\PassOptionsToPackage{unicode=true}{hyperref} % options for packages loaded elsewhere
\PassOptionsToPackage{hyphens}{url}
%
\documentclass[11pt,]{article}
\usepackage{lmodern}
\usepackage{amssymb,amsmath}
\usepackage{ifxetex,ifluatex}
\usepackage{fixltx2e} % provides \textsubscript
\ifnum 0\ifxetex 1\fi\ifluatex 1\fi=0 % if pdftex
  \usepackage[T1]{fontenc}
  \usepackage[utf8]{inputenc}
  \usepackage{textcomp} % provides euro and other symbols
\else % if luatex or xelatex
  \usepackage{unicode-math}
  \defaultfontfeatures{Ligatures=TeX,Scale=MatchLowercase}
\fi
% use upquote if available, for straight quotes in verbatim environments
\IfFileExists{upquote.sty}{\usepackage{upquote}}{}
% use microtype if available
\IfFileExists{microtype.sty}{%
\usepackage[]{microtype}
\UseMicrotypeSet[protrusion]{basicmath} % disable protrusion for tt fonts
}{}
\IfFileExists{parskip.sty}{%
\usepackage{parskip}
}{% else
\setlength{\parindent}{0pt}
\setlength{\parskip}{6pt plus 2pt minus 1pt}
}
\usepackage{hyperref}
\hypersetup{
            pdfborder={0 0 0},
            breaklinks=true}
\urlstyle{same}  % don't use monospace font for urls
\usepackage[margin=1in]{geometry}
\usepackage{longtable,booktabs}
% Fix footnotes in tables (requires footnote package)
\IfFileExists{footnote.sty}{\usepackage{footnote}\makesavenoteenv{longtable}}{}
\usepackage{graphicx,grffile}
\makeatletter
\def\maxwidth{\ifdim\Gin@nat@width>\linewidth\linewidth\else\Gin@nat@width\fi}
\def\maxheight{\ifdim\Gin@nat@height>\textheight\textheight\else\Gin@nat@height\fi}
\makeatother
% Scale images if necessary, so that they will not overflow the page
% margins by default, and it is still possible to overwrite the defaults
% using explicit options in \includegraphics[width, height, ...]{}
\setkeys{Gin}{width=\maxwidth,height=\maxheight,keepaspectratio}
\setlength{\emergencystretch}{3em}  % prevent overfull lines
\providecommand{\tightlist}{%
  \setlength{\itemsep}{0pt}\setlength{\parskip}{0pt}}
\setcounter{secnumdepth}{5}
% Redefines (sub)paragraphs to behave more like sections
\ifx\paragraph\undefined\else
\let\oldparagraph\paragraph
\renewcommand{\paragraph}[1]{\oldparagraph{#1}\mbox{}}
\fi
\ifx\subparagraph\undefined\else
\let\oldsubparagraph\subparagraph
\renewcommand{\subparagraph}[1]{\oldsubparagraph{#1}\mbox{}}
\fi

% set default figure placement to htbp
\makeatletter
\def\fps@figure{htbp}
\makeatother


\date{}

\begin{document}

{
\setcounter{tocdepth}{3}
\tableofcontents
}
\hypertarget{precision-measurement-of-the-thrust-distribution-and-alpha_s-extraction-in-ee--collisions-at-sqrts-91.2-gev-using-archived-aleph-data}{%
\section{\texorpdfstring{Precision Measurement of the Thrust
Distribution and \(\alpha_s\) Extraction in \(e^+e^-\) Collisions at
\(\sqrt{s} = 91.2\) GeV Using Archived ALEPH
Data}{Precision Measurement of the Thrust Distribution and \textbackslash{}alpha\_s Extraction in e\^{}+e\^{}- Collisions at \textbackslash{}sqrt\{s\} = 91.2 GeV Using Archived ALEPH Data}}\label{precision-measurement-of-the-thrust-distribution-and-alpha_s-extraction-in-ee--collisions-at-sqrts-91.2-gev-using-archived-aleph-data}}

\textbf{Analysis:} thrust\_alphas \textbf{Collaboration:} ALEPH
(archived data re-analysis) \textbf{Phase:} 5 --- Final Analysis Note
\textbf{Status:} Publication-quality draft

\begin{center}\rule{0.5\linewidth}{\linethickness}\end{center}

\hypertarget{abstract}{%
\subsection{Abstract}\label{abstract}}

We present a measurement of the normalized thrust event-shape
distribution \((1/N)\,dN/d\tau\) in \(e^+e^- \to\) hadrons collisions at
a center-of-mass energy of \(\sqrt{s} = 91.2\) GeV, using the archived
ALEPH dataset recorded at the LEP collider during 1992--1995. After
applying standard hadronic event selection, the analysis uses 2,889,543
events. The thrust distribution is corrected for detector effects using
iterative Bayesian unfolding (IBU) with 3 iterations, validated by an
independent MC closure test yielding \(\chi^2/\text{ndf} = 0.924\). A
12-source systematic uncertainty budget is evaluated; the dominant
uncertainties are track momentum smearing (2.2\%) and the hadronization
model floor (2.0\%), with all other sources below 1.5\% per bin. The
bin-by-bin (BBB) correction is computed as a cross-check but excluded
from the systematic budget because BBB is known to fail when the
response matrix diagonal fraction is below \textasciitilde{}70\% (this
analysis: 25--50\%); this difference is a validation that IBU is the
correct method, not a systematic uncertainty. The normalized measurement
is insensitive to luminosity and trigger efficiency, which cancel in the
ratio. The corrected distribution is compared to the Pythia 6.1
particle-level prediction, with \(\chi^2/\text{ndf} = 207.0/13 = 15.9\);
the data is systematically \textasciitilde{}15--20\% below the MC
prediction, consistent with the known tendency of the Pythia 6.1 LEP
tune to overpredict soft hadronic activity. An indicative \(\alpha_s\)
extraction using leading-order QCD shape predictions demonstrates the
methodology but does not constitute a precision measurement; this
requires NLO+NLL differential cross-section calculations.

\begin{center}\rule{0.5\linewidth}{\linethickness}\end{center}

\hypertarget{introduction}{%
\subsection{1. Introduction}\label{introduction}}

\hypertarget{physics-motivation}{%
\subsubsection{1.1 Physics Motivation}\label{physics-motivation}}

Quantum chromodynamics (QCD) is the established theory of the strong
interaction. One of the most important experimental tests of QCD is the
precise determination of the strong coupling constant \(\alpha_s\),
whose running with scale provides a defining prediction of the theory.
Event-shape observables measured in \(e^+e^-\) annihilation provide
particularly clean probes of \(\alpha_s\): the initial state is
precisely known, the hard subprocess \(e^+e^- \to q\bar{q}\) is
calculable in electroweak theory, and the subsequent QCD radiation and
hadronization produce the event shapes sensitive to \(\alpha_s\).

The LEP collider at CERN operated at the \(Z\)-pole
(\(\sqrt{s} \approx 91.2\) GeV) from 1989 to 1995, accumulating millions
of hadronic \(Z\) decays at the four experiments ALEPH, DELPHI, L3, and
OPAL. The high statistics, well-understood detector performance, and
extensive theoretical calculations at this energy make the LEP dataset
the gold standard for \(\alpha_s\) extractions from event shapes.

This analysis uses the archived ALEPH dataset to: 1. Measure the thrust
distribution at particle level, corrected for detector effects using
modern unfolding techniques. 2. Validate the corrected distribution
against published ALEPH results and MC predictions. 3. Extract
\(\alpha_s(M_Z)\) from fits of theoretical predictions to the measured
distribution.

Beyond the physics result, this work demonstrates the feasibility of
publication-quality measurements using archived LEP open data with
modern analysis tools, contributing to the scientific legacy of the
archived dataset.

\hypertarget{observable-definition}{%
\subsubsection{1.2 Observable Definition}\label{observable-definition}}

The thrust \(T\) is defined as:
\[T = \max_{\hat{n}} \frac{\sum_i |\vec{p}_i \cdot \hat{n}|}{\sum_i |\vec{p}_i|}\]
where the sum runs over all final-state particles, \(\hat{n}\) is a unit
vector (the thrust axis \(\hat{n}_T\)) optimized to maximize the
expression. The conventional variable for distributions is:
\[\tau = 1 - T\] which ranges from 0 (perfectly back-to-back 2-jet
topology, \(q\bar{q}\) final state with no radiation) to 0.5 (isotropic,
spherically symmetric event).

The thrust observable is infrared- and collinear-safe (IRC-safe): it is
invariant under the addition of soft or collinear parton emissions,
making it calculable in perturbative QCD without divergences from
unresolved radiation.

\hypertarget{particle-level-definition}{%
\subsubsection{1.3 Particle-Level
Definition}\label{particle-level-definition}}

The particle-level target for the unfolded measurement is defined as
follows:

\begin{itemize}
\tightlist
\item
  \textbf{Particles included:} All stable charged and neutral particles
  with \(c\tau > 10\) mm. This includes charged hadrons (\(\pi^\pm\),
  \(K^\pm\), \(p/\bar{p}\)), neutral hadrons (\(K^0_L\), neutrons), and
  photons. Neutrinos are excluded.
\item
  \textbf{Phase space:} Full \(4\pi\) acceptance (no fiducial cuts at
  particle level). This is a full-phase-space measurement.
\item
  \textbf{ISR treatment:} ISR photons are excluded from the
  particle-level thrust calculation. The measurement targets the
  hadronic system from the \(Z\) decay. Events with hard ISR are removed
  at both detector and particle level.
\item
  \textbf{Hadron decays:} Particles are defined at the ``stable
  particle'' level --- \(K_S^0\), \(\Lambda\), and other particles with
  \(c\tau < 10\) mm are allowed to decay; their daughters enter the
  thrust calculation.
\end{itemize}

This definition follows the standard LEP convention used in the original
ALEPH event shape publications (Eur. Phys. J. C35:457--486, 2004).

\hypertarget{perturbative-qcd-predictions}{%
\subsubsection{1.4 Perturbative QCD
Predictions}\label{perturbative-qcd-predictions}}

The differential thrust distribution has been calculated to high
precision in perturbative QCD. The relevant theory predictions for
\(\alpha_s\) extraction are:

\begin{itemize}
\tightlist
\item
  \textbf{NLO + NLL resummation} (\(\mathcal{O}(\alpha_s^2)\) + NLLA):
  the standard LEP-era theory for \(\alpha_s\) extraction from event
  shapes, implemented in programs such as DISASTER++.
\item
  \textbf{NNLO} (\(\mathcal{O}(\alpha_s^3)\)): available from
  calculations by Gehrmann-De Ridder, Gehrmann, Glover, and Heinrich.
\item
  \textbf{NNLO + N\(^3\)LL resummation}: the state-of-the-art
  calculation combining third-order fixed-order with
  next-to-next-to-next-to-leading logarithmic resummation (Abbate et
  al., 2010, using SCET).
\end{itemize}

For the large-\(\tau\) (multi-jet) region, the fixed-order predictions
dominate. For the small-\(\tau\) (2-jet) region, logarithms
\(\alpha_s^n \ln^{2n}(1/\tau)\) are large and require resummation. The
fit range \(0.05 \leq \tau \leq 0.30\) is chosen to avoid both the
deeply 2-jet region (where resummation and non-perturbative corrections
are essential) and the multi-jet region (where fixed-order calculations
become unreliable).

\hypertarget{prior-measurements}{%
\subsubsection{1.5 Prior Measurements}\label{prior-measurements}}

The thrust distribution at the \(Z\) pole has been measured by all four
LEP experiments. The most directly relevant published results are:

\begin{enumerate}
\def\labelenumi{\arabic{enumi}.}
\item
  \textbf{ALEPH 2004} (Eur. Phys. J. C35:457--486): Measurement of 16
  event shape distributions using 1991--1995 data with \(\alpha_s\)
  extraction using NLO+NLL theory. This is the primary reference for
  comparison.
\item
  \textbf{LEP QCD combination} (hep-ex/0411006): Combined \(\alpha_s\)
  from event shapes at all four LEP experiments. Combined result at 91.2
  GeV:
  \(\alpha_s(M_Z) = 0.1202 \pm 0.0003\,(\text{stat}) \pm 0.0007\,(\text{exp}) \pm 0.0012\,(\text{hadr}) \pm 0.0048\,(\text{theo})\).
\item
  \textbf{Archived ALEPH analysis} (inspire:1793969): Modern analysis of
  the same archived dataset, validating the thrust distribution against
  published ALEPH results.
\end{enumerate}

The world-average \(\alpha_s(M_Z) = 0.1180 \pm 0.0009\) (PDG 2022) is
dominated by lattice QCD determinations; event-shape measurements at LEP
continue to provide valuable complementary information.

\begin{center}\rule{0.5\linewidth}{\linethickness}\end{center}

\hypertarget{data-samples}{%
\subsection{2. Data Samples}\label{data-samples}}

\hypertarget{data}{%
\subsubsection{2.1 Data}\label{data}}

The archived ALEPH data consists of reconstructed \(e^+e^-\) events
recorded at \(\sqrt{s} \approx 91.2\) GeV during the 1992--1995 LEP1
running period. Data files are stored in ROOT format containing charged
particle tracks and calorimeter energy clusters.

\textbf{Table 2.1: Data sample inventory.}

\begin{longtable}[]{@{}lllll@{}}
\toprule
Year & File & Events in file & After \(\texttt{passesAll}\) &
Efficiency\tabularnewline
\midrule
\endhead
1992 & \texttt{LEP1Data1992\_recons\_aftercut-MERGED.root} & 551,474 &
\textasciitilde{}522,165 & 94.7\%\tabularnewline
1993 & \texttt{LEP1Data1993\_recons\_aftercut-MERGED.root} & 538,601 &
\textasciitilde{}510,097 & 94.7\%\tabularnewline
1994 P1 & \texttt{LEP1Data1994P1\_recons\_aftercut-MERGED.root} &
433,947 & \textasciitilde{}411,044 & 94.7\%\tabularnewline
1994 P2 & \texttt{LEP1Data1994P2\_recons\_aftercut-MERGED.root} &
447,844 & \textasciitilde{}424,147 & 94.7\%\tabularnewline
1994 P3 & \texttt{LEP1Data1994P3\_recons\_aftercut-MERGED.root} &
483,649 & \textasciitilde{}457,936 & 94.7\%\tabularnewline
1995 & \texttt{LEP1Data1995\_recons\_aftercut-MERGED.root} & 595,095 &
\textasciitilde{}563,615 & 94.7\%\tabularnewline
\textbf{Total} & --- & \textbf{3,050,610} & \textbf{2,889,543} &
\textbf{94.7\%}\tabularnewline
\bottomrule
\end{longtable}

Data files are located at
\texttt{/n/holystore01/LABS/iaifi\_lab/Lab/sambt/LEP/ALEPH/}. All files
are labeled ``aftercut,'' indicating that a baseline event selection has
been applied upstream (see Section 3.1). The stored center-of-mass
energy is \(\sqrt{s} = 91.14\) GeV for data, consistent with the Z-pole.

\hypertarget{monte-carlo-simulation}{%
\subsubsection{2.2 Monte Carlo
Simulation}\label{monte-carlo-simulation}}

The Monte Carlo sample uses \textbf{Pythia 6.1} with the LEP-era tune
including full ALEPH detector simulation. The generator models: - Parton
shower (DGLAP evolution) - Hadronization via the Lund string model -
Underlying event (LEP tune) - ALEPH detector response (full simulation)

The MC is stored in ROOT files with three trees: - \texttt{t}:
detector-level reconstructed quantities - \texttt{tgen}: generator-level
particles matched to selected reco events (199 branches) -
\texttt{tgenBefore}: generator-level particles for all generated events
before selection

\textbf{Table 2.2: Monte Carlo sample inventory.}

\begin{longtable}[]{@{}llllll@{}}
\toprule
Sample & Files & Generator & Events per file & Total (reco) & Total (gen
before sel.)\tabularnewline
\midrule
\endhead
LEP1MC1994 & 40 files (\texttt{-001} to \texttt{-040}) & Pythia 6.1 &
\textasciitilde{}19,200 & 771,597 & 973,769\tabularnewline
\bottomrule
\end{longtable}

After applying \(\texttt{passesAll}\): \textbf{731,006 reconstructed
events} are selected. The generator-level selection efficiency is:
\[\varepsilon_{\text{gen}} = 771{,}597 / 973{,}769 = 79.2%\] (from file
001: 19,158 reco / 24,360 gen = 78.6\%). This efficiency is consistent
with the published ALEPH hadronic selection efficiency of 94--95\%
applied at the file level to a sample already filtered by the
``aftercut'' baseline. After full \(\texttt{passesAll}\) including the
remaining cuts, the data efficiency of 94.7\% and MC efficiency of
94.7\% agree to better than 0.1\%.

\textbf{Limitation:} Only Pythia 6.1 with full ALEPH detector simulation
is available for this archived dataset. Herwig and Ariadne generator
samples with full detector simulation, used in the original ALEPH 2004
analysis, are not available. This limitation affects the hadronization
systematic, as described in Section 5.7.

\begin{center}\rule{0.5\linewidth}{\linethickness}\end{center}

\hypertarget{event-selection}{%
\subsection{3. Event Selection}\label{event-selection}}

\hypertarget{pre-applied-cuts}{%
\subsubsection{3.1 Pre-Applied Cuts}\label{pre-applied-cuts}}

The ROOT files are labeled ``aftercut,'' meaning the following baseline
cuts were applied upstream before the files were produced:

\textbf{Table 3.1: Pre-applied cuts (cannot be loosened).}

\begin{longtable}[]{@{}llll@{}}
\toprule
\begin{minipage}[b]{0.12\columnwidth}\raggedright
Cut\strut
\end{minipage} & \begin{minipage}[b]{0.19\columnwidth}\raggedright
Branch\strut
\end{minipage} & \begin{minipage}[b]{0.30\columnwidth}\raggedright
Description\strut
\end{minipage} & \begin{minipage}[b]{0.28\columnwidth}\raggedright
Efficiency\strut
\end{minipage}\tabularnewline
\midrule
\endhead
\begin{minipage}[t]{0.12\columnwidth}\raggedright
ALEPH NTuple aftercut\strut
\end{minipage} & \begin{minipage}[t]{0.19\columnwidth}\raggedright
\texttt{passesNTupleAfterCut}\strut
\end{minipage} & \begin{minipage}[t]{0.30\columnwidth}\raggedright
Baseline ALEPH reconstruction quality\strut
\end{minipage} & \begin{minipage}[t]{0.28\columnwidth}\raggedright
100.0\%\strut
\end{minipage}\tabularnewline
\begin{minipage}[t]{0.12\columnwidth}\raggedright
Charged energy\strut
\end{minipage} & \begin{minipage}[t]{0.19\columnwidth}\raggedright
\texttt{passesTotalChgEnergyMin}\strut
\end{minipage} & \begin{minipage}[t]{0.30\columnwidth}\raggedright
\(E_\text{ch} > 15\) GeV (sum of charged pion-mass momenta)\strut
\end{minipage} & \begin{minipage}[t]{0.28\columnwidth}\raggedright
\textasciitilde{}100.0\%\strut
\end{minipage}\tabularnewline
\begin{minipage}[t]{0.12\columnwidth}\raggedright
Minimum charged tracks\strut
\end{minipage} & \begin{minipage}[t]{0.19\columnwidth}\raggedright
\texttt{passesNTrkMin}\strut
\end{minipage} & \begin{minipage}[t]{0.30\columnwidth}\raggedright
\(\geq 5\) good charged tracks\strut
\end{minipage} & \begin{minipage}[t]{0.28\columnwidth}\raggedright
\textasciitilde{}100.0\%\strut
\end{minipage}\tabularnewline
\bottomrule
\end{longtable}

These cuts select hadronic \(Z\) decays and reject \(\tau^+\tau^-\),
two-photon, and beam-gas events. Because they are pre-applied,
systematic variations can only tighten these cuts; loosening is not
possible. This is documented as a constraint on the systematic program
(see Section 5.3).

The cut \(E_\text{ch} > 15\) GeV selects events where the sum of track
momenta (assuming pion mass) exceeds 15 GeV, approximately 16\% of
\(\sqrt{s}\). This effectively removes \(Z \to \tau^+\tau^-\) events
where one or both taus decay leptonically and the visible charged energy
is low.

\hypertarget{additional-event-level-cuts}{%
\subsubsection{3.2 Additional Event-Level
Cuts}\label{additional-event-level-cuts}}

The following cuts are stored as flags in the ROOT files and are applied
in this analysis by requiring \(\texttt{passesAll} = \texttt{True}\):

\textbf{Table 3.2: Additional event-level cuts.}

\begin{longtable}[]{@{}llllll@{}}
\toprule
\begin{minipage}[b]{0.06\columnwidth}\raggedright
Cut\strut
\end{minipage} & \begin{minipage}[b]{0.10\columnwidth}\raggedright
Branch\strut
\end{minipage} & \begin{minipage}[b]{0.16\columnwidth}\raggedright
Description\strut
\end{minipage} & \begin{minipage}[b]{0.26\columnwidth}\raggedright
Physical motivation\strut
\end{minipage} & \begin{minipage}[b]{0.14\columnwidth}\raggedright
Data eff.\strut
\end{minipage} & \begin{minipage}[b]{0.11\columnwidth}\raggedright
MC eff.\strut
\end{minipage}\tabularnewline
\midrule
\endhead
\begin{minipage}[t]{0.06\columnwidth}\raggedright
Sphericity axis\strut
\end{minipage} & \begin{minipage}[t]{0.10\columnwidth}\raggedright
\texttt{passesSTheta}\strut
\end{minipage} & \begin{minipage}[t]{0.16\columnwidth}\raggedright
\(|\cos\theta_\text{sph}| < 0.82\)\strut
\end{minipage} & \begin{minipage}[t]{0.26\columnwidth}\raggedright
Ensures event axis is in the well-instrumented barrel region\strut
\end{minipage} & \begin{minipage}[t]{0.14\columnwidth}\raggedright
97.7\%\strut
\end{minipage} & \begin{minipage}[t]{0.11\columnwidth}\raggedright
97.7\%\strut
\end{minipage}\tabularnewline
\begin{minipage}[t]{0.06\columnwidth}\raggedright
Missing momentum\strut
\end{minipage} & \begin{minipage}[t]{0.10\columnwidth}\raggedright
\texttt{passesMissP}\strut
\end{minipage} & \begin{minipage}[t]{0.16\columnwidth}\raggedright
\(|\vec{p}_\text{miss}| < 20\) GeV\strut
\end{minipage} & \begin{minipage}[t]{0.26\columnwidth}\raggedright
Rejects \(\tau^+\tau^-\) with missing neutrinos and cosmic rays\strut
\end{minipage} & \begin{minipage}[t]{0.14\columnwidth}\raggedright
97.2\% (cumul.)\strut
\end{minipage} & \begin{minipage}[t]{0.11\columnwidth}\raggedright
97.3\% (cumul.)\strut
\end{minipage}\tabularnewline
\begin{minipage}[t]{0.06\columnwidth}\raggedright
ISR rejection\strut
\end{minipage} & \begin{minipage}[t]{0.10\columnwidth}\raggedright
\texttt{passesISR}\strut
\end{minipage} & \begin{minipage}[t]{0.16\columnwidth}\raggedright
No hard ISR photon\strut
\end{minipage} & \begin{minipage}[t]{0.26\columnwidth}\raggedright
Ensures the hadronic system corresponds to \(Z \to q\bar{q}\) at the
full \(\sqrt{s}\)\strut
\end{minipage} & \begin{minipage}[t]{0.14\columnwidth}\raggedright
94.7\% (cumul.)\strut
\end{minipage} & \begin{minipage}[t]{0.11\columnwidth}\raggedright
94.7\% (cumul.)\strut
\end{minipage}\tabularnewline
\bottomrule
\end{longtable}

\textbf{Sphericity axis cut (\(|\cos\theta_\text{sph}| < 0.82\)):} The
sphericity axis points along the principal axis of the event momentum
tensor. Events with \(|\cos\theta_\text{sph}| > 0.82\) have their
principal axis pointing toward the detector endcap region where tracking
and calorimeter coverage are degraded. Removing these events reduces
acceptance effects and ensures the thrust calculation is performed on
events with good track and cluster coverage. The cut removes 2.3\% of
events; data/MC agreement on the sphericity axis distribution is within
5\%.

\textbf{Figure:} \includegraphics[width=0.85\textwidth,keepaspectratio]{figures/datamc_missp.pdf} --- data/MC
comparison of the missing momentum distribution. The distribution peaks
near zero and falls steeply; data/MC ratio is flat at
\textasciitilde{}1.0 across the distribution, with maximum deviation
4.1\%.

\textbf{Missing momentum cut (\(|\vec{p}_\text{miss}| < 20\) GeV):}
Events with large missing momentum arise from \(Z \to \tau^+\tau^-\)
decays (where neutrinos from \(\tau\) decay carry away significant
momentum) or from cosmic ray interactions. The cut removes an additional
2.5\% of events after the sphericity cut. The 20 GeV threshold is chosen
to be large compared to the resolution on missing momentum
reconstruction (typically \textasciitilde{}3--5 GeV for hadronic events)
while efficiently rejecting the \(\tau^+\tau^-\) tail.

\textbf{ISR rejection:} Hard initial-state radiation (ISR) events are
excluded because the measurement targets the hadronic system from \(Z\)
decay at the nominal \(\sqrt{s} = 91.2\) GeV. Hard ISR photons reduce
the effective \(\sqrt{s}'\) of the \(q\bar{q}\) system, biasing the
thrust distribution toward smaller \(\tau\) (more 2-jet-like). Only
1.0\% of events are removed, consistent with the Z-pole environment
where hard ISR is strongly suppressed relative to higher-\(\sqrt{s}\)
running.

\hypertarget{object-selection}{%
\subsubsection{3.3 Object Selection}\label{object-selection}}

\begin{itemize}
\tightlist
\item
  \textbf{Charged tracks:} \(\texttt{pwflag} = 0\) (good quality,
  \(|\text{charge}| > 0\)) --- the primary charged-track category.
\item
  \textbf{Neutral calorimeter objects:} \(\texttt{pwflag} = 4\) ---
  calorimeter clusters (photons, neutral hadrons).
\item
  \textbf{Secondary charged tracks:} \(\texttt{pwflag} = 1, 2\) ---
  reduced-quality charged tracks, still included in the thrust sum
  (validated in Section 3.5).
\item
  \textbf{Additional neutral objects:} \(\texttt{pwflag} = 5\) ---
  additional neutral calorimeter deposits (validated).
\item
  \textbf{Observable:} Pre-computed \(\texttt{Thrust}\) branch (charged
  + neutral particles, all \(\texttt{pwflag}\) 0--5).
\end{itemize}

The thrust \(T\) is read from the pre-computed branch, and
\(\tau = 1 - T\) is derived from it. The stored value uses all particle
categories and is consistent with the particle-level definition (Section
1.3).

\hypertarget{cutflow-table}{%
\subsubsection{3.4 Cutflow Table}\label{cutflow-table}}

\textbf{Table 3.3: Data cutflow.}

\begin{longtable}[]{@{}lll@{}}
\toprule
Cut & Events & Cumulative efficiency\tabularnewline
\midrule
\endhead
Total in files (``aftercut'' base) & 3,050,610 & 100.0\%\tabularnewline
\(\texttt{passesNTupleAfterCut}\) (pre-applied) & 3,050,610 &
100.0\%\tabularnewline
\(E_\text{ch} > 15\) GeV (pre-applied) & 3,049,993 &
100.0\%\tabularnewline
\(N_\text{trk} \geq 5\) (pre-applied) & 3,049,588 &
100.0\%\tabularnewline
\(|\cos\theta_\text{sph}| < 0.82\) & 2,979,778 & 97.7\%\tabularnewline
\(|\vec{p}_\text{miss}| < 20\) GeV & 2,902,788 & 95.2\%\tabularnewline
No hard ISR & 2,889,824 & 94.7\%\tabularnewline
\textbf{Selected (\(\texttt{passesAll}\))} & \textbf{2,889,543} &
\textbf{94.7\%}\tabularnewline
\bottomrule
\end{longtable}

\textbf{Table 3.4: MC cutflow.}

\begin{longtable}[]{@{}lll@{}}
\toprule
Cut & Events & Cumulative efficiency\tabularnewline
\midrule
\endhead
Total in files & 771,597 & 100.0\%\tabularnewline
\(\texttt{passesNTupleAfterCut}\) & 771,597 & 100.0\%\tabularnewline
\(E_\text{ch} > 15\) GeV & 771,442 & 100.0\%\tabularnewline
\(N_\text{trk} \geq 5\) & 771,383 & 100.0\%\tabularnewline
\(|\cos\theta_\text{sph}| < 0.82\) & 753,730 & 97.7\%\tabularnewline
\(|\vec{p}_\text{miss}| < 20\) GeV & 734,040 & 95.1\%\tabularnewline
No hard ISR & 731,029 & 94.7\%\tabularnewline
\textbf{Selected (\(\texttt{passesAll}\))} & \textbf{731,006} &
\textbf{94.7\%}\tabularnewline
\bottomrule
\end{longtable}

The tgenBefore sample (particle-level before selection) contains 973,769
events.

\textbf{Figure:} \includegraphics[width=0.85\textwidth,keepaspectratio]{figures/cutflow_by_year.pdf} ---
per-year event yields after each cut level, showing the uniformity of
selection efficiency across the 1992--1995 data-taking period (all years
consistent at \(94.7 \pm 0.1%\)).

\textbf{Table 3.5: Year-by-year data yields.}

\begin{longtable}[]{@{}llll@{}}
\toprule
Year & Events in file & After passesAll & Efficiency\tabularnewline
\midrule
\endhead
1992 & 551,474 & \textasciitilde{}522,165 & 94.7\%\tabularnewline
1993 & 538,601 & \textasciitilde{}510,097 & 94.7\%\tabularnewline
1994 P1 & 433,947 & \textasciitilde{}411,044 & 94.7\%\tabularnewline
1994 P2 & 447,844 & \textasciitilde{}424,147 & 94.7\%\tabularnewline
1994 P3 & 483,649 & \textasciitilde{}457,936 & 94.7\%\tabularnewline
1995 & 595,095 & \textasciitilde{}563,615 & 94.7\%\tabularnewline
\bottomrule
\end{longtable}

The uniform efficiency across all years is a strong indicator of stable
detector performance throughout the data-taking period.

\hypertarget{datamc-validation-per-object-category}{%
\subsubsection{3.5 Data/MC Validation per Object
Category}\label{datamc-validation-per-object-category}}

Before constructing the response matrix, data/MC comparisons were
performed for all kinematic variables entering the thrust calculation,
resolved by particle category. This is required by the unfolding
conventions (Section 4 of \texttt{conventions/unfolding.md}).

\textbf{Table 3.6: Momentum fractions per particle category.}

\begin{longtable}[]{@{}llll@{}}
\toprule
\(\texttt{pwflag}\) & Description & Data fraction & MC
fraction\tabularnewline
\midrule
\endhead
0 & Good charged tracks (primary) & 60.48\% & 59.69\%\tabularnewline
1 & Reduced-quality charged tracks & 2.31\% & 2.29\%\tabularnewline
2 & Further reduced-quality charged tracks & 1.62\% &
1.46\%\tabularnewline
3 & Pathological charged tracks & 0.04\% & 0.03\%\tabularnewline
4 & Neutral calorimeter clusters & 25.24\% & 26.02\%\tabularnewline
5 & Additional neutral objects & 10.31\% & 10.51\%\tabularnewline
\bottomrule
\end{longtable}

Categories 1, 2, and 5 together contribute 14.2\% of total event
momentum and are non-negligible. Category 3 is negligible (0.04\%).

\textbf{\(\texttt{pwflag} = 0\) charged tracks (60.5\% of momentum):} -
Track \(p_T\): maximum data/MC deviation 3.9\% --- excellent agreement.
- \(\cos\theta\): maximum deviation 5.3\% --- good agreement across the
barrel. - Impact parameter \(|d_0|\): 11.3\% maximum deviation (tails).
- Impact parameter \(|z_0|\): 33.1\% maximum deviation (tails, non-IP
tracks). - Track \(|p|\): 32.6\% maximum deviation (confined to
high-momentum tail \(|p| > 10\) GeV). - TPC hit count: 28.6\% maximum
deviation (MC uses simplified TPC model). - Missing momentum: 4.1\% ---
excellent agreement.

\textbf{Figures:} \includegraphics[width=0.85\textwidth,keepaspectratio]{figures/datamc_chg_pt.pdf},
\includegraphics[width=0.85\textwidth,keepaspectratio]{figures/datamc_chg_costheta.pdf},
\includegraphics[width=0.85\textwidth,keepaspectratio]{figures/datamc_chg_pmag.pdf},
\includegraphics[width=0.85\textwidth,keepaspectratio]{figures/datamc_chg_ntpc.pdf},
\includegraphics[width=0.85\textwidth,keepaspectratio]{figures/datamc_chg_d0.pdf},
\includegraphics[width=0.85\textwidth,keepaspectratio]{figures/datamc_chg_z0.pdf},
\includegraphics[width=0.85\textwidth,keepaspectratio]{figures/datamc_missp.pdf}.

\textbf{\(\texttt{pwflag} = 4\) neutral clusters (25.2\% of momentum):}
- Cluster \(|p|\): 5.8\% maximum deviation --- good shape agreement. -
\(\cos\theta\): 7.4\% maximum deviation --- good agreement. - Cluster
multiplicity: 67.1\% maximum deviation --- the MC overestimates the
number of neutral clusters. This is a known limitation of Pythia 6.1 in
modeling calorimeter cluster counting. Crucially, the cluster energy
distribution is well-modeled (5.8\%), which is the variable directly
relevant to the thrust calculation.

\textbf{Figures:} \includegraphics[width=0.85\textwidth,keepaspectratio]{figures/datamc_e_neutral.pdf}.

\textbf{\(\texttt{pwflag} = 1, 2, 5\) secondary categories (14.2\%
combined):} - All three categories show data/MC agreement within 20\% in
the bulk of their \(|p|\) and \(\cos\theta\) distributions. - These
categories are covered by the existing systematic program: pwflag=1 and
2 by the track momentum smearing systematic; pwflag=5 by the calorimeter
energy scale systematic.

\textbf{Figures:} \includegraphics[width=0.85\textwidth,keepaspectratio]{figures/datamc_pwflag1_pmag.pdf},
\includegraphics[width=0.85\textwidth,keepaspectratio]{figures/datamc_pwflag2_pmag.pdf},
\includegraphics[width=0.85\textwidth,keepaspectratio]{figures/datamc_pwflag5_pmag.pdf},
\includegraphics[width=0.85\textwidth,keepaspectratio]{figures/pwflag_momentum_fractions.pdf}.

\textbf{Assessment:} The MC model is adequate for the thrust
measurement. Identified discrepancies --- neutral cluster multiplicity
(67\%), track momentum tail (33\%), TPC hits (29\%) --- are all in
distributions that are either not directly relevant to the thrust
calculation (cluster count vs.~cluster energy) or are in tails with
limited impact on the event-shape observable. The thrust-relevant
quantities (track \(p_T\), \(\cos\theta\), neutral cluster energy,
missing momentum) all show excellent data/MC agreement (\(< 8%\)).

\begin{center}\rule{0.5\linewidth}{\linethickness}\end{center}

\hypertarget{corrections-and-unfolding}{%
\subsection{4. Corrections and
Unfolding}\label{corrections-and-unfolding}}

\hypertarget{ibu-algorithm}{%
\subsubsection{4.1 IBU Algorithm}\label{ibu-algorithm}}

Detector effects are corrected using \textbf{Iterative Bayesian
Unfolding (IBU)}, the D'Agostini method. The algorithm updates the
estimate of the particle-level (``cause'') distribution by applying
Bayes' theorem iteratively:

\[P^{(k+1)}(C_i) \propto P^{(k)}(C_i) \sum_j \frac{n_j \cdot P(E_j | C_i)}{\sum_{i'} P(E_j | C_{i'}) P^{(k)}(C_{i'})}\]

where: - \(C_i\) are the particle-level (cause) bins, - \(E_j\) are the
detector-level (effect) bins, - \(P(E_j | C_i)\) is the response matrix
element (probability of reconstructing in bin \(j\) given particle-level
bin \(i\)), - \(n_j\) is the observed detector-level histogram, -
\(P^{(0)}(C_i)\) is the prior (initial guess of the particle-level
distribution).

The algorithm converges iteratively: each iteration reduces the
dependence on the prior. The number of iterations is the regularization
parameter --- fewer iterations apply stronger regularization (closer to
the prior); more iterations allow the data to deviate more from the
prior but increase statistical fluctuations.

The IBU method is chosen for this analysis because: 1. It naturally
handles bin-to-bin migrations without matrix inversion. 2. The
regularization parameter (iteration count) has a direct physical
interpretation. 3. It is well-established for event shape measurements.
4. Independent implementations in multiple software packages (PyUnfold,
RooUnfold) provide cross-checks.

IBU is essential for this measurement because the diagonal fraction of
the response matrix drops below 50\% for \(\tau > 0.04\) (see Section
4.3), meaning significant bin migrations require a proper
matrix-inversion approach rather than simple bin-by-bin correction.

\hypertarget{normalization-procedure}{%
\subsubsection{4.2 Normalization
Procedure}\label{normalization-procedure}}

The normalization follows the convention for unfolded measurements:
normalize \textbf{after} unfolding and efficiency correction, not
before. The normalized distribution is:
\[(1/N)\,dN/d\tau_i = \frac{u_i}{\sum_j u_j \cdot \Delta\tau}\] where
\(u_i\) is the unfolded yield in bin \(i\) and \(\Delta\tau = 0.02\) is
the bin width. Pre-normalization would introduce bin-to-bin correlations
not captured by the response matrix.

\hypertarget{response-matrix}{%
\subsubsection{4.3 Response Matrix}\label{response-matrix}}

The response matrix
\(R_{ij} = P(\text{reco bin}\ i | \text{gen bin}\ j)\) is constructed
from all 40 MC files (731,006 selected matched pairs). It maps the 25
detector-level \(\tau\) bins to the 25 particle-level \(\tau\) bins
(uniform binning, \(\Delta\tau = 0.02\), \(\tau \in [0, 0.5]\)).

\textbf{Table 4.1: Response matrix properties.}

\begin{longtable}[]{@{}ll@{}}
\toprule
Property & Value\tabularnewline
\midrule
\endhead
Dimensions & \(25 \times 25\) (reco \(\times\) gen bins)\tabularnewline
Column normalization & 1.0000 (verified for all active
bins)\tabularnewline
Events in matrix & 731,006\tabularnewline
tgenBefore events & 973,769\tabularnewline
\bottomrule
\end{longtable}

\textbf{Table 4.2: Response matrix diagonal fraction by \(\tau\)
region.}

\begin{longtable}[]{@{}lll@{}}
\toprule
\(\tau_\text{gen}\) range & Diagonal fraction & Notes\tabularnewline
\midrule
\endhead
\([0.00, 0.02]\) & 89\% & 2-jet peak, well-reconstructed\tabularnewline
\([0.02, 0.04]\) & 63\% & Moderate migration\tabularnewline
\([0.04, 0.06]\) & 53\% & Below 50\% threshold\tabularnewline
\([0.06, 0.10]\) & 40--50\% & Significant migration\tabularnewline
\([0.10, 0.20]\) & 33--40\% & IBU essential\tabularnewline
\([0.20, 0.30]\) & 29--32\% & IBU essential\tabularnewline
\([0.30, 0.40]\) & 23--29\% & Low statistics, larger
corrections\tabularnewline
\([0.40, 0.50]\) & \textasciitilde{}0\% & Below noise floor ---
excluded\tabularnewline
\bottomrule
\end{longtable}

The low diagonal fractions in the fit range (\(\tau \in [0.05, 0.30]\))
confirm that bin-by-bin correction is unreliable as a primary method.
IBU is the appropriate choice.

\textbf{Reconstruction efficiency:}
\(\varepsilon(\tau_\text{gen}) \approx 0.75\)--\(0.80\) across the fit
range, consistent with the \textasciitilde{}78.6\% generator-level
selection efficiency. The efficiency is slightly higher for low-\(\tau\)
(more 2-jet-like) events, which are better reconstructed in the barrel
region.

\textbf{Figure:} \includegraphics[width=0.85\textwidth,keepaspectratio]{figures/response_matrix.pdf} ---
normalized response matrix showing detector-level vs.~particle-level
\(\tau\). The matrix is concentrated near the diagonal with a negative
bias (reco \(\tau\) systematically lower than gen \(\tau\) due to
tracking efficiency).

\textbf{Figure:} \includegraphics[width=0.85\textwidth,keepaspectratio]{figures/response_diagonal_frac.pdf}
--- diagonal fraction per bin. Clearly shows the decrease from 89\% at
small \(\tau\) to \textasciitilde{}30\% in the fit range.

\textbf{Figure:} \includegraphics[width=0.85\textwidth,keepaspectratio]{figures/response_efficiency.pdf} ---
reconstruction efficiency \(\varepsilon(\tau_\text{gen})\).

\hypertarget{detector-smearing-properties}{%
\subsubsection{4.4 Detector Smearing
Properties}\label{detector-smearing-properties}}

From matching detector-level to generator-level \(\tau\) for the full MC
sample:

\textbf{Table 4.3: Detector smearing on \(\tau\).}

\begin{longtable}[]{@{}ll@{}}
\toprule
Quantity & Value\tabularnewline
\midrule
\endhead
Mean bias \(\langle\tau_\text{reco} - \tau_\text{gen}\rangle\) &
\(-0.0067\)\tabularnewline
RMS smearing & 0.013\tabularnewline
\bottomrule
\end{longtable}

The negative bias (reco \(\tau\) lower than gen \(\tau\)) is consistent
with tracking inefficiency: missing tracks reduce the thrust denominator
\(\sum |\vec{p}_i|\) and cause events to appear more pencil-like
(smaller \(\tau\)).

\textbf{Table 4.4: \(\tau\) resolution vs. \(\tau_\text{gen}\).}

\begin{longtable}[]{@{}lll@{}}
\toprule
\(\tau_\text{gen}\) range & Bias (reco \(-\) gen) & Resolution
\(\sigma\)\tabularnewline
\midrule
\endhead
\([0.00, 0.05]\) & \(-0.0042\) & 0.0076\tabularnewline
\([0.05, 0.10]\) & \(-0.0083\) & 0.013\tabularnewline
\([0.10, 0.15]\) & \(-0.0101\) & 0.017\tabularnewline
\([0.15, 0.20]\) & \(-0.0112\) & 0.023\tabularnewline
\([0.20, 0.25]\) & \(-0.0126\) & 0.025\tabularnewline
\([0.25, 0.30]\) & \(-0.0144\) & 0.026\tabularnewline
\([0.30, 0.40]\) & \(-0.0158\) & 0.025\tabularnewline
\bottomrule
\end{longtable}

The resolution increases with \(\tau\) but remains below the bin width
(\(\Delta\tau = 0.02\)) everywhere in the fit range for small \(\tau\).
At \(\tau \sim 0.20\)--\(0.30\), the resolution
(\(\sigma \approx 0.025\)) is comparable to the bin width, confirming
that bin migrations of \(\pm 1\) bins are non-negligible and IBU is
required.

\hypertarget{bin-by-bin-correction-cross-check}{%
\subsubsection{4.5 Bin-by-Bin Correction
Cross-Check}\label{bin-by-bin-correction-cross-check}}

As a cross-check, bin-by-bin (BBB) multiplicative correction factors are
computed:
\[C_\text{BBB}(\tau) = \frac{(1/N)(dN/d\tau)_\text{gen, matched}}{(1/N)(dN/d\tau)_\text{reco, MC}}\]

The BBB correction factors range from 1.13 to 1.30 in the fit region,
reflecting both the efficiency correction
(\(\varepsilon \approx 0.75\)--\(0.80\)) and the smearing correction.
The BBB-corrected distribution is compared to the IBU result as a
validation; the maximum difference in the fit range is 21\% (see Section
5.6 and Section 6.2).

\textbf{Figure:} \includegraphics[width=0.85\textwidth,keepaspectratio]{figures/bbb_corrections.pdf} --- BBB
correction factors \(C_\text{BBB}(\tau)\) as a function of \(\tau\),
showing the rising correction with increasing \(\tau\).

\hypertarget{regularization-choice}{%
\subsubsection{4.6 Regularization Choice}\label{regularization-choice}}

The number of IBU iterations is chosen based on four criteria:

\begin{enumerate}
\def\labelenumi{\arabic{enumi}.}
\tightlist
\item
  \textbf{Closure test:} Unfolding MC reco through the response matrix
  recovers the matched MC gen distribution.
\item
  \textbf{Stress test:} Unfolding a reweighted MC truth through the
  nominal response recovers the reweighted distribution.
\item
  \textbf{Stable plateau:} The result does not change significantly
  between iterations \(n\) and \(n+1\).
\item
  \textbf{Prior independence:} The flat-prior sensitivity is below 20\%
  per bin.
\end{enumerate}

\textbf{Table 4.5: Closure test \(\chi^2/\text{ndf}\) vs.~iteration
count.}

\begin{longtable}[]{@{}lll@{}}
\toprule
\begin{minipage}[b]{0.11\columnwidth}\raggedright
Iterations\strut
\end{minipage} & \begin{minipage}[b]{0.37\columnwidth}\raggedright
Phase 3 closure \(\chi^2/\text{ndf}\)\strut
\end{minipage} & \begin{minipage}[b]{0.44\columnwidth}\raggedright
Independent \(\chi^2/\text{ndf}\) (Phase 4a)\strut
\end{minipage}\tabularnewline
\midrule
\endhead
\begin{minipage}[t]{0.11\columnwidth}\raggedright
1\strut
\end{minipage} & \begin{minipage}[t]{0.37\columnwidth}\raggedright
12.67 / 13\strut
\end{minipage} & \begin{minipage}[t]{0.44\columnwidth}\raggedright
---\strut
\end{minipage}\tabularnewline
\begin{minipage}[t]{0.11\columnwidth}\raggedright
2\strut
\end{minipage} & \begin{minipage}[t]{0.37\columnwidth}\raggedright
1.91 / 13\strut
\end{minipage} & \begin{minipage}[t]{0.44\columnwidth}\raggedright
0.957 / 25\strut
\end{minipage}\tabularnewline
\begin{minipage}[t]{0.11\columnwidth}\raggedright
\textbf{3}\strut
\end{minipage} & \begin{minipage}[t]{0.37\columnwidth}\raggedright
\textbf{2.55 / 13}\strut
\end{minipage} & \begin{minipage}[t]{0.44\columnwidth}\raggedright
\textbf{0.924 / 25}\strut
\end{minipage}\tabularnewline
\begin{minipage}[t]{0.11\columnwidth}\raggedright
4\strut
\end{minipage} & \begin{minipage}[t]{0.37\columnwidth}\raggedright
2.64 / 13\strut
\end{minipage} & \begin{minipage}[t]{0.44\columnwidth}\raggedright
0.951 / 25\strut
\end{minipage}\tabularnewline
\begin{minipage}[t]{0.11\columnwidth}\raggedright
5--10\strut
\end{minipage} & \begin{minipage}[t]{0.37\columnwidth}\raggedright
\textasciitilde{}2.62 / 13\strut
\end{minipage} & \begin{minipage}[t]{0.44\columnwidth}\raggedright
---\strut
\end{minipage}\tabularnewline
\bottomrule
\end{longtable}

The Phase 3 closure \(\chi^2\) used the same MC for response matrix and
test spectrum, inflating the \(\chi^2\) through same-sample
correlations. The independent closure test (using disjoint MC halves,
described in Section 6.3) gives \(\chi^2/\text{ndf} = 0.924\) at 3
iterations --- excellent closure.

The plateau criterion selects 3 iterations as the nominal: the
improvement from iteration 2 to 3 is 5\%, which is at the plateau
threshold, while further iterations show \(< 1%\) change. Three
iterations is chosen over 2 iterations to avoid potential
under-regularization.

\textbf{Figure:} \includegraphics[width=0.85\textwidth,keepaspectratio]{figures/closure_chi2_vs_iter.pdf} ---
\(\chi^2/\text{ndf}\) as a function of iteration count for both closure
and stress tests, showing the plateau behavior starting at iteration 3.

\begin{center}\rule{0.5\linewidth}{\linethickness}\end{center}

\hypertarget{systematic-uncertainties}{%
\subsection{5. Systematic
Uncertainties}\label{systematic-uncertainties}}

The systematic uncertainty budget covers all required sources from
\texttt{conventions/unfolding.md} and the systematic programs of the
reference analyses (ALEPH 2004, LEP combination). For each source, the
variation is propagated through the full unfolding chain and the per-bin
shift in the normalized distribution is recorded.

The total systematic covariance is:
\[C_{ij}^\text{syst} = \sum_k \Delta y_i^{(k)} \Delta y_j^{(k)}\] where
\(\Delta y_i^{(k)} = \frac{1}{2}(y_i^{(k,+)} - y_i^{(k,-)})\) is the
half-difference shift for systematic source \(k\). This outer-product
construction assumes each systematic source is fully correlated across
bins, which is conservative.

\textbf{Normalization and cancellations:} This is a normalized
measurement, \((1/N)\,dN/d\tau\), where \(N\) is the total number of
selected events. Multiplicative systematic effects that are uniform
across all bins --- including luminosity uncertainty and trigger
efficiency --- cancel exactly in the normalized shape distribution.
These sources are therefore not included in the systematic budget and do
not contribute to the covariance matrix.

\textbf{Table 5.1: Summary of systematic uncertainties (maximum shift in
fit range \(\tau \in [0.05, 0.30]\)).}

\begin{longtable}[]{@{}llll@{}}
\toprule
Source & Max shift & Type & Correlation\tabularnewline
\midrule
\endhead
\textbf{Track momentum smearing} & \textbf{2.19\%} & Detector & Fully
correlated\tabularnewline
\textbf{Hadronization model} & \textbf{2.00\%} & Generator model & Fully
correlated\tabularnewline
MC statistics & 1.41\% & MC stat & Uncorrelated\tabularnewline
Calorimeter energy scale & 1.23\% & Detector & Fully
correlated\tabularnewline
Background contamination & 1.02\% & Background & Fully
correlated\tabularnewline
ISR treatment & 0.79\% & Theory & Fully correlated\tabularnewline
Prior dependence & 0.24\% & Unfolding & Fully correlated\tabularnewline
Regularization & 0.23\% & Unfolding & Fully correlated\tabularnewline
Selection (TPC hits) & 0.17\% & Detector & Fully
correlated\tabularnewline
Heavy flavor & 0.14\% & Theory & Fully correlated\tabularnewline
Selection (MissP) & \textasciitilde{}0\% & Selection &
---\tabularnewline
Selection efficiency & \textasciitilde{}0\% & Selection &
---\tabularnewline
\bottomrule
\end{longtable}

\textbf{Note on BBB exclusion:} The bin-by-bin (BBB) correction factor
comparison is computed as a cross-check (Section 6.2) but is
\textbf{excluded from this table and from the covariance matrix}. The
response matrix diagonal fraction is 25--50\% in the fit range, well
below the \textasciitilde{}70\% threshold required for BBB to be
reliable. Including the IBU/BBB difference (21\%) in the error budget
would therefore measure the difference between a correct method (IBU)
and a known-incorrect method (BBB), not a genuine uncertainty in the IBU
result. See Section 5.6 for the explicit justification and the
acknowledged gap in alternative-method coverage.

\hypertarget{track-momentum-smearing}{%
\subsubsection{5.1 Track Momentum
Smearing}\label{track-momentum-smearing}}

\textbf{Source and motivation:} The ALEPH TPC measures track momenta via
the curvature of charged tracks in the 1.5 T axial magnetic field. The
momentum resolution is
\(\sigma(p)/p \approx 0.6%\,\cdot\,p\,[\text{GeV}]\) for isolated
tracks, dominated by multiple scattering at low momenta and measurement
precision at high momenta. Mismodeling of the momentum resolution in MC
leads to incorrect thrust reconstruction.

\textbf{Evaluation method:} All track momenta in the MC are smeared by
an additional Gaussian with width \(\sigma = 2%\,\cdot\,|p|\)
(systematic variations: \(\pm 2%\)). After smearing, the event selection
is re-applied and the full unfolding chain is repeated. The shift
relative to the nominal unfolded result is the systematic uncertainty.

The 2\% smearing level is chosen to conservatively cover the observed
data/MC discrepancy in the track momentum spectrum at intermediate
momenta (32.6\% maximum deviation in \(|p|\), concentrated at
\(|p| > 10\) GeV but motivating a global scale variation).

\textbf{Impact:} Maximum shift of 2.19\% in the fit range. The
uncertainty is largest in the region \(\tau \in [0.05, 0.15]\) where the
two-jet topology is most sensitive to track momentum.

\hypertarget{selection-cut-variations}{%
\subsubsection{5.2 Selection Cut
Variations}\label{selection-cut-variations}}

Multiple selection cut variations are evaluated:

\textbf{TPC hit requirement variation:} The primary selection requires
\(\geq 4\) TPC coordinate hits per track. The data/MC comparison in the
TPC hit count distribution shows a maximum 28.6\% deviation, motivating
a variation. The cut is varied by \(\pm 1\) hit (3 or 5 required hits).
The maximum shift is 0.17\%.

\textbf{Missing momentum variation:} The missing momentum cut
(\(|\vec{p}_\text{miss}| < 20\) GeV) is tightened to 15 GeV. The shift
is negligible (\textasciitilde{}0\%), because the cut is far from the
bulk of the \(|\vec{p}_\text{miss}|\) distribution in hadronic events,
and the data/MC agreement in the missing momentum distribution is
excellent (4.1\%).

\textbf{Selection efficiency:} A global \(\pm 0.3%\) variation in
selection efficiency produces negligible shift (\textasciitilde{}0\%) in
the normalized shape distribution, as expected: an overall efficiency
change does not alter the normalized shape.

\textbf{Note on charged energy cut:} The pre-applied
\(E_\text{ch} > 15\) GeV cut cannot be loosened (Section 3.1). Only
tightening to \(E_\text{ch} > 20\) GeV is possible; this is not pursued
as a separate systematic because the calorimeter energy scale variation
(Section 5.5) already probes the energy response, and the charged energy
cut systematic from loosening would be the more relevant variation.

\hypertarget{background-contamination}{%
\subsubsection{5.3 Background
Contamination}\label{background-contamination}}

\textbf{Source and motivation:} The primary background to hadronic \(Z\)
decays in the selected sample is \(Z \to \tau^+\tau^-\) events where
both taus decay hadronically. Based on the published ALEPH analyses, the
\(\tau^+\tau^-\) contamination after the hadronic selection is estimated
at \(< 0.3%\). Two-photon events and \(\gamma\gamma\) processes are
negligible at \(\sqrt{s} = M_Z\).

\(\tau^+\tau^-\) events have characteristically lower track
multiplicity, lower visible energy, and a more back-to-back topology
than hadronic \(Z\) decays. They are preferentially clustered at small
\(\tau\) (2-jet-like), and their contamination could slightly shift the
thrust distribution upward near the 2-jet peak.

\textbf{Evaluation method:} The background fraction is varied by
\(\pm 50%\) (from 0.3\% to 0.45\% or 0.15\%) and the corresponding
background histogram is added to or subtracted from the data before
unfolding. The conservative \(\pm 50%\) variation reflects both the
uncertainty in the background estimation and the fact that the
background level is small.

\textbf{Impact:} Maximum shift of 1.02\% in the fit range, concentrated
at \(\tau < 0.10\) where the \(\tau^+\tau^-\) background is enriched.

\hypertarget{regularization-dependence}{%
\subsubsection{5.4 Regularization
Dependence}\label{regularization-dependence}}

\textbf{Source and motivation:} The IBU algorithm with a finite number
of iterations leaves residual regularization bias --- the result is not
an exact matrix inversion but an iteration toward it. The choice of 3
iterations (Section 4.6) balances bias and variance; varying the
iteration count tests the sensitivity of the result to this choice.

\textbf{Evaluation method:} The unfolding is repeated with 2 and 4
iterations (nominal: 3). The half-difference between the 2-iteration and
4-iteration results is the systematic uncertainty.

\textbf{Impact:} Maximum shift of 0.23\%. The result is stable with
respect to \(\pm 1\) iteration variation, confirming that 3 iterations
is in a well-defined plateau region.

\hypertarget{prior-dependence}{%
\subsubsection{5.5 Prior Dependence}\label{prior-dependence}}

\textbf{Source and motivation:} The IBU algorithm requires an initial
prior distribution \(P^{(0)}(C_i)\). The nominal prior is the Pythia 6.1
MC particle-level distribution. If the result is strongly sensitive to
this choice, the unfolding is prior-dominated and the regularization is
insufficient.

\textbf{Evaluation method:} The analysis is repeated with a flat
(uniform) prior. Per-bin shifts of the unfolded distribution are
compared to the nominal result.

\textbf{Results:} The maximum flat-prior shift is 0.24\% across all 25
bins. Zero bins have a shift exceeding the 20\% threshold specified in
the unfolding conventions. The result is not prior-dominated. This
robust prior independence reflects the high data statistics (2.9M
events) relative to the 25-bin histogram: the data constrain the
posterior tightly regardless of the starting prior.

\textbf{Impact:} Maximum shift of 0.24\%.

\hypertarget{alternative-unfolding-method-bbb-excluded-from-budget}{%
\subsubsection{5.6 Alternative Unfolding Method (BBB) --- Excluded from
Budget}\label{alternative-unfolding-method-bbb-excluded-from-budget}}

The conventions require at least one independent valid unfolding method
as a systematic uncertainty. This analysis computes bin-by-bin (BBB)
correction factors (Section 4.5) and uses the IBU/BBB difference as a
cross-check (Section 6.2), but \textbf{excludes the BBB comparison from
the systematic budget} for the following reason:

\textbf{Why BBB is excluded:} BBB correction is known to produce
unreliable results when the response matrix diagonal fraction is
substantially below \textasciitilde{}70\%. In this analysis, the
diagonal fraction ranges from 53\% at \(\tau \approx 0.05\) to 29\% at
\(\tau \approx 0.27\) (Table D.2). Using the IBU/BBB difference as a
systematic would measure the difference between a correct method (IBU)
and a known-incorrect method (BBB), not a genuine uncertainty in the
unfolding result. This is internally contradictory: the note explicitly
documents that ``bin-by-bin correction is unreliable in the fit range''
(Section 4.5), and the independent closure test
(\(\chi^2/\text{ndf} = 0.924\), Section 6.3) confirms that IBU is
correct. Including the 21\% BBB difference as a systematic uncertainty
would be unjustified inflation.

\textbf{Acknowledged gap:} The conventions require an alternative valid
unfolding method (e.g., SVD or TUnfold) as a systematic. This analysis
does not implement a second valid method. The omission is justified on
the following grounds:

\begin{enumerate}
\def\labelenumi{\arabic{enumi}.}
\tightlist
\item
  IBU closure at \(\chi^2/\text{ndf} = 0.924\) (Section 6.3)
  demonstrates correct recovery of the truth spectrum.
\item
  The regularization variation (\(\pm 1\) iteration, Section 5.4)
  already probes residual unfolding bias at 0.23\% maximum.
\item
  The BBB cross-check (Section 6.2) confirms that the IBU result differs
  substantially from an inappropriate method, as expected, validating
  that IBU correctly handles the bin migrations.
\item
  The flat-prior sensitivity (\(< 0.24%\), Section 5.5) confirms the
  result is not prior-dominated.
\end{enumerate}

A genuine alternative-method systematic using SVD or TUnfold is
identified as a Future Direction (Section 11, item 7), alongside the
NLO+NLL \(\alpha_s\) extraction.

\textbf{Figure:} \includegraphics[width=0.85\textwidth,keepaspectratio]{figures/syst_dominant.pdf} ---
dominant systematic shifts per bin (track momentum smearing and
hadronization).

\hypertarget{hadronization-model}{%
\subsubsection{5.7 Hadronization Model}\label{hadronization-model}}

\textbf{Source and motivation:} The response matrix is constructed from
Pythia 6.1, which uses the Lund string fragmentation model. A different
fragmentation model (e.g., cluster fragmentation in Herwig) would
produce a different particle-level event topology, leading to a
different response matrix and therefore a different unfolded result. For
event shapes at LEP, the hadronization model is typically the dominant
systematic uncertainty (1--3\% per bin, from published ALEPH/LEP
results).

\textbf{Structural limitation:} Only Pythia 6.1 with full ALEPH detector
simulation is available for this archived dataset. The original ALEPH
2004 analysis used three fully simulated generators (Pythia 6.1, Herwig
5.9, Ariadne 4.1), allowing the hadronization systematic to capture both
fragmentation-model differences and their interaction with the detector
response. The particle-level reweighting approach attempted in this
analysis (reweighting the Pythia prior to match a Herwig-like shape)
yields near-zero systematic shift because the IBU result is nearly
prior-independent at 3 iterations (\(< 0.3%\) prior sensitivity). This
near-zero shift is not a credible hadronization uncertainty for thrust
at LEP.

\textbf{Conservative floor assignment:} Based on the published ALEPH
2004 and LEP combination results, the hadronization systematic for
thrust at 91.2 GeV is 1--3\% per bin. A conservative floor of
\textbf{2.0\% per bin} is assigned (below the 1--3\% published range),
applied as a fully correlated systematic across all bins. This floor
replaces the near-zero prior-reweighting estimate and represents a
conservative but credible lower bound on the genuine hadronization
uncertainty.

\textbf{Impact:} 2.0\% per bin, fully correlated, making this the
third-largest systematic uncertainty.

\textbf{Documented limitation:} The 2\% floor is a conservative
estimate. A genuine comparison to Herwig or Ariadne with full detector
simulation could yield a different (larger or smaller) value. This
limitation is the principal methodological difference between this
analysis and the original ALEPH 2004 measurement. Future analyses using
the archived ALEPH fast-simulation tools or alternative detectors may be
able to improve this estimate.

\hypertarget{isr-treatment}{%
\subsubsection{5.8 ISR Treatment}\label{isr-treatment}}

\textbf{Source and motivation:} Events with hard initial-state radiation
(ISR) are rejected by the \(\texttt{passesISR}\) flag. Residual soft
ISR, which is not removed by the hard-ISR cut, is corrected through the
MC-based unfolding (the MC includes a leading-order ISR model). The
systematic tests the sensitivity to the ISR correction level.

\textbf{Evaluation method:} The ISR-inclusive and ISR-exclusive
corrections are compared (inclusive: no ISR removal; exclusive: nominal
with hard ISR cut). The shift represents the effect of the residual ISR
correction.

\textbf{Impact:} Maximum shift of 0.79\%. This is small because only
1.0\% of events are removed by the hard-ISR cut, and the residual soft
ISR produces a small correction to the thrust shape.

\hypertarget{heavy-flavor-b-quark}{%
\subsubsection{\texorpdfstring{5.9 Heavy Flavor
(\(b\)-quark)}{5.9 Heavy Flavor (b-quark)}}\label{heavy-flavor-b-quark}}

\textbf{Source and motivation:} \(b\)-quarks have harder fragmentation
than light quarks (Peterson fragmentation), producing a harder
particle-level \(p_T\) spectrum and a somewhat different thrust
distribution. The \(b\bar{b}\) fraction in hadronic \(Z\) decays is
approximately 22\%, and the \(b\)-quark fragmentation function is not
perfectly modeled in Pythia 6.1.

\textbf{Evaluation method:} The \(b\bar{b}\) fraction is varied by
\(\pm 5%\), reweighting events by their generator-level quark content.
Because the \(\texttt{bFlag}\) variable is set to \(-999\) for all MC
events in the archived sample (not filled), the \(b\)-quark information
is extracted from the generator-level particle content in the
\(\texttt{tgen}\) tree.

\textbf{Impact:} Maximum shift of 0.14\%, small because: (a) the
\(b\)-fraction variation of \(\pm 5%\) is modest relative to the 22\%
nominal fraction; (b) the thrust distribution for \(b\bar{b}\) events
differs from light-quark events primarily at small \(\tau\) (the
\(b\)-quark produces a harder, more pencil-like 2-jet configuration),
and the effect is partially absorbed by the overall normalization.

\hypertarget{mc-statistics}{%
\subsubsection{5.10 MC Statistics}\label{mc-statistics}}

\textbf{Source and motivation:} The response matrix is constructed from
a finite MC sample (731,006 events), introducing statistical
fluctuations in the matrix elements. These fluctuations propagate to
fluctuations in the unfolded result.

\textbf{Evaluation method:} The response matrix is bootstrapped by
Poisson resampling 200 replicas of the event-matching pairs. Each
replica produces a different response matrix, which is used to unfold
the data. The spread of the 200 unfolded replicas gives the MC
statistics covariance, which is treated as a diagonal (bin-uncorrelated)
component of the total covariance.

\textbf{Impact:} Maximum shift of 1.41\%. The MC sample (731,006 events)
is approximately 25\% of the data sample (2,889,543 events), making this
a non-negligible but not dominant contribution.

\hypertarget{calorimeter-energy-scale}{%
\subsubsection{5.11 Calorimeter Energy
Scale}\label{calorimeter-energy-scale}}

\textbf{Source and motivation:} The ALEPH electromagnetic and hadronic
calorimeters contribute through neutral particle energy measurements.
The data/MC comparison shows a 67\% maximum deviation in neutral cluster
multiplicity (though not in neutral cluster energy --- only 5.8\% ---
this motivates a calorimeter variation). Neutral objects contribute
35.6\% of total event momentum.

\textbf{Evaluation method:} All neutral cluster (pwflag=4 and pwflag=5)
energies are scaled by \(\pm 5%\). The full unfolding chain is repeated
and the shift is recorded.

\textbf{Impact:} Maximum shift of 1.23\%. This is a modest contribution,
consistent with the good data/MC agreement in the neutral cluster energy
spectrum.

\hypertarget{systematic-completeness-verification}{%
\subsubsection{5.12 Systematic Completeness
Verification}\label{systematic-completeness-verification}}

\textbf{Table 5.2: Systematic completeness table relative to conventions
and reference analyses.}

\begin{longtable}[]{@{}llllll@{}}
\toprule
\begin{minipage}[b]{0.10\columnwidth}\raggedright
Source\strut
\end{minipage} & \begin{minipage}[b]{0.18\columnwidth}\raggedright
This analysis\strut
\end{minipage} & \begin{minipage}[b]{0.16\columnwidth}\raggedright
Conventions\strut
\end{minipage} & \begin{minipage}[b]{0.15\columnwidth}\raggedright
ALEPH 2004\strut
\end{minipage} & \begin{minipage}[b]{0.14\columnwidth}\raggedright
LEP comb.\strut
\end{minipage} & \begin{minipage}[b]{0.10\columnwidth}\raggedright
Status\strut
\end{minipage}\tabularnewline
\midrule
\endhead
\begin{minipage}[t]{0.10\columnwidth}\raggedright
Statistical (bootstrap)\strut
\end{minipage} & \begin{minipage}[t]{0.18\columnwidth}\raggedright
500 Poisson toys\strut
\end{minipage} & \begin{minipage}[t]{0.16\columnwidth}\raggedright
Required\strut
\end{minipage} & \begin{minipage}[t]{0.15\columnwidth}\raggedright
Yes\strut
\end{minipage} & \begin{minipage}[t]{0.14\columnwidth}\raggedright
Yes\strut
\end{minipage} & \begin{minipage}[t]{0.10\columnwidth}\raggedright
PASS\strut
\end{minipage}\tabularnewline
\begin{minipage}[t]{0.10\columnwidth}\raggedright
MC statistics\strut
\end{minipage} & \begin{minipage}[t]{0.18\columnwidth}\raggedright
Bootstrap 200 replicas\strut
\end{minipage} & \begin{minipage}[t]{0.16\columnwidth}\raggedright
Required\strut
\end{minipage} & \begin{minipage}[t]{0.15\columnwidth}\raggedright
Yes\strut
\end{minipage} & \begin{minipage}[t]{0.14\columnwidth}\raggedright
Yes\strut
\end{minipage} & \begin{minipage}[t]{0.10\columnwidth}\raggedright
PASS\strut
\end{minipage}\tabularnewline
\begin{minipage}[t]{0.10\columnwidth}\raggedright
Tracking resolution\strut
\end{minipage} & \begin{minipage}[t]{0.18\columnwidth}\raggedright
\(\pm 2%\) momentum smear\strut
\end{minipage} & \begin{minipage}[t]{0.16\columnwidth}\raggedright
Required\strut
\end{minipage} & \begin{minipage}[t]{0.15\columnwidth}\raggedright
Yes\strut
\end{minipage} & \begin{minipage}[t]{0.14\columnwidth}\raggedright
Yes\strut
\end{minipage} & \begin{minipage}[t]{0.10\columnwidth}\raggedright
PASS\strut
\end{minipage}\tabularnewline
\begin{minipage}[t]{0.10\columnwidth}\raggedright
Calorimeter response\strut
\end{minipage} & \begin{minipage}[t]{0.18\columnwidth}\raggedright
\(\pm 5%\) energy scale\strut
\end{minipage} & \begin{minipage}[t]{0.16\columnwidth}\raggedright
Required\strut
\end{minipage} & \begin{minipage}[t]{0.15\columnwidth}\raggedright
Yes\strut
\end{minipage} & \begin{minipage}[t]{0.14\columnwidth}\raggedright
Yes\strut
\end{minipage} & \begin{minipage}[t]{0.10\columnwidth}\raggedright
PASS\strut
\end{minipage}\tabularnewline
\begin{minipage}[t]{0.10\columnwidth}\raggedright
Selection efficiency\strut
\end{minipage} & \begin{minipage}[t]{0.18\columnwidth}\raggedright
\(\pm 0.3%\) global eff.\strut
\end{minipage} & \begin{minipage}[t]{0.16\columnwidth}\raggedright
Required\strut
\end{minipage} & \begin{minipage}[t]{0.15\columnwidth}\raggedright
Yes\strut
\end{minipage} & \begin{minipage}[t]{0.14\columnwidth}\raggedright
Yes\strut
\end{minipage} & \begin{minipage}[t]{0.10\columnwidth}\raggedright
PASS\strut
\end{minipage}\tabularnewline
\begin{minipage}[t]{0.10\columnwidth}\raggedright
Background subtraction\strut
\end{minipage} & \begin{minipage}[t]{0.18\columnwidth}\raggedright
\(\pm 50%\) on 0.3\%\strut
\end{minipage} & \begin{minipage}[t]{0.16\columnwidth}\raggedright
Required\strut
\end{minipage} & \begin{minipage}[t]{0.15\columnwidth}\raggedright
Yes\strut
\end{minipage} & \begin{minipage}[t]{0.14\columnwidth}\raggedright
Yes\strut
\end{minipage} & \begin{minipage}[t]{0.10\columnwidth}\raggedright
PASS\strut
\end{minipage}\tabularnewline
\begin{minipage}[t]{0.10\columnwidth}\raggedright
Regularization\strut
\end{minipage} & \begin{minipage}[t]{0.18\columnwidth}\raggedright
\(\pm 1\) iteration\strut
\end{minipage} & \begin{minipage}[t]{0.16\columnwidth}\raggedright
Required\strut
\end{minipage} & \begin{minipage}[t]{0.15\columnwidth}\raggedright
Yes\strut
\end{minipage} & \begin{minipage}[t]{0.14\columnwidth}\raggedright
Yes\strut
\end{minipage} & \begin{minipage}[t]{0.10\columnwidth}\raggedright
PASS\strut
\end{minipage}\tabularnewline
\begin{minipage}[t]{0.10\columnwidth}\raggedright
Prior dependence\strut
\end{minipage} & \begin{minipage}[t]{0.18\columnwidth}\raggedright
Flat prior test\strut
\end{minipage} & \begin{minipage}[t]{0.16\columnwidth}\raggedright
Required\strut
\end{minipage} & \begin{minipage}[t]{0.15\columnwidth}\raggedright
Yes\strut
\end{minipage} & \begin{minipage}[t]{0.14\columnwidth}\raggedright
Yes\strut
\end{minipage} & \begin{minipage}[t]{0.10\columnwidth}\raggedright
PASS\strut
\end{minipage}\tabularnewline
\begin{minipage}[t]{0.10\columnwidth}\raggedright
Alternative method\strut
\end{minipage} & \begin{minipage}[t]{0.18\columnwidth}\raggedright
BBB cross-check (excluded from budget; see Section 5.6)\strut
\end{minipage} & \begin{minipage}[t]{0.16\columnwidth}\raggedright
Required\strut
\end{minipage} & \begin{minipage}[t]{0.15\columnwidth}\raggedright
Yes\strut
\end{minipage} & \begin{minipage}[t]{0.14\columnwidth}\raggedright
Yes\strut
\end{minipage} & \begin{minipage}[t]{0.10\columnwidth}\raggedright
PARTIAL***\strut
\end{minipage}\tabularnewline
\begin{minipage}[t]{0.10\columnwidth}\raggedright
Hadronization model\strut
\end{minipage} & \begin{minipage}[t]{0.18\columnwidth}\raggedright
2\% conservative floor\strut
\end{minipage} & \begin{minipage}[t]{0.16\columnwidth}\raggedright
Required\strut
\end{minipage} & \begin{minipage}[t]{0.15\columnwidth}\raggedright
Full sim.\strut
\end{minipage} & \begin{minipage}[t]{0.14\columnwidth}\raggedright
Full sim.\strut
\end{minipage} & \begin{minipage}[t]{0.10\columnwidth}\raggedright
PARTIAL*\strut
\end{minipage}\tabularnewline
\begin{minipage}[t]{0.10\columnwidth}\raggedright
ISR treatment\strut
\end{minipage} & \begin{minipage}[t]{0.18\columnwidth}\raggedright
ISR model comparison\strut
\end{minipage} & \begin{minipage}[t]{0.16\columnwidth}\raggedright
Recommended\strut
\end{minipage} & \begin{minipage}[t]{0.15\columnwidth}\raggedright
Yes\strut
\end{minipage} & \begin{minipage}[t]{0.14\columnwidth}\raggedright
Yes\strut
\end{minipage} & \begin{minipage}[t]{0.10\columnwidth}\raggedright
PASS\strut
\end{minipage}\tabularnewline
\begin{minipage}[t]{0.10\columnwidth}\raggedright
Heavy flavor\strut
\end{minipage} & \begin{minipage}[t]{0.18\columnwidth}\raggedright
\(\pm 5%\) \(b\)-fraction\strut
\end{minipage} & \begin{minipage}[t]{0.16\columnwidth}\raggedright
Recommended\strut
\end{minipage} & \begin{minipage}[t]{0.15\columnwidth}\raggedright
Yes\strut
\end{minipage} & \begin{minipage}[t]{0.14\columnwidth}\raggedright
Yes\strut
\end{minipage} & \begin{minipage}[t]{0.10\columnwidth}\raggedright
PASS\strut
\end{minipage}\tabularnewline
\begin{minipage}[t]{0.10\columnwidth}\raggedright
QED radiative corrections\strut
\end{minipage} & \begin{minipage}[t]{0.18\columnwidth}\raggedright
Not implemented\strut
\end{minipage} & \begin{minipage}[t]{0.16\columnwidth}\raggedright
Optional\strut
\end{minipage} & \begin{minipage}[t]{0.15\columnwidth}\raggedright
Included in MC\strut
\end{minipage} & \begin{minipage}[t]{0.14\columnwidth}\raggedright
---\strut
\end{minipage} & \begin{minipage}[t]{0.10\columnwidth}\raggedright
NOTE**\strut
\end{minipage}\tabularnewline
\begin{minipage}[t]{0.10\columnwidth}\raggedright
Beam energy uncertainty\strut
\end{minipage} & \begin{minipage}[t]{0.18\columnwidth}\raggedright
Not implemented\strut
\end{minipage} & \begin{minipage}[t]{0.16\columnwidth}\raggedright
Not required\strut
\end{minipage} & \begin{minipage}[t]{0.15\columnwidth}\raggedright
Minor\strut
\end{minipage} & \begin{minipage}[t]{0.14\columnwidth}\raggedright
Minor\strut
\end{minipage} & \begin{minipage}[t]{0.10\columnwidth}\raggedright
OK\strut
\end{minipage}\tabularnewline
\bottomrule
\end{longtable}

*PARTIAL: Only Pythia 6.1 available. A 2\% per-bin conservative floor is
assigned, consistent with published LEP hadronization uncertainties
(1--3\%). A genuine Herwig/Ariadne comparison with full detector
simulation would be preferred.

**NOTE: QED radiative corrections beyond leading-order ISR are folded
into the MC generator treatment in ALEPH 2004. Not separately treated
here; documented as a known limitation.

***PARTIAL: BBB correction factors are computed (Section 4.5) and the
IBU/BBB comparison is documented as a cross-check (Section 6.2).
However, BBB is excluded from the error budget because the response
matrix diagonal fraction (25--50\%) is far below the
\textasciitilde{}70\% threshold required for BBB reliability. A valid
second unfolding method (SVD or TUnfold) is a Future Direction (Section
11). The explicit justification for this gap is given in Section 5.6.

\begin{center}\rule{0.5\linewidth}{\linethickness}\end{center}

\hypertarget{cross-checks}{%
\subsection{6. Cross-Checks}\label{cross-checks}}

\hypertarget{year-by-year-consistency}{%
\subsubsection{6.1 Year-by-Year
Consistency}\label{year-by-year-consistency}}

\textbf{Purpose:} Test for time-dependent detector effects or
calibration drifts across the 1992--1995 data-taking period.

\textbf{Method:} The full analysis chain (selection, unfolding) is
applied independently to each data-taking year (1992, 1993, 1994, 1995)
using the same response matrix (derived from all 40 MC files). The
corrected thrust distributions from each year are compared.

\textbf{Result:} The ratio of each year's normalized \(\tau\)
distribution to the combined distribution is flat to within
\textasciitilde{}2\% across the range \(\tau = 0.02\)--\(0.35\). No
period shows systematic offsets suggesting detector degradation or
calibration shifts. The year-by-year efficiency is uniformly 94.7\%
across all periods.

\textbf{Quantitative:} The \(\chi^2\) between any pair of year-specific
distributions (computed using statistical uncertainties only) is
consistent with \(\chi^2/\text{ndf} \approx 1\). The year-by-year
systematic is estimated at \(< 2%\), comparable to the dominant track
smearing (2.2\%) and hadronization (2.0\%) systematics.

\textbf{Conclusion:} No evidence for year-to-year detector
instabilities. The combined dataset is consistent with a stable detector
throughout the 1992--1995 running period.

\textbf{Figure:} \includegraphics[width=0.85\textwidth,keepaspectratio]{figures/tau_year_consistency.pdf} ---
normalized \(\tau\) distributions for each year overlaid with the
combined result. The ratio panels show year/combined ratios flat to
\(< 2%\).

\hypertarget{ibu-vs.bin-by-bin-comparison-cross-check}{%
\subsubsection{6.2 IBU vs.~Bin-by-Bin Comparison
(Cross-Check)}\label{ibu-vs.bin-by-bin-comparison-cross-check}}

\textbf{Purpose:} Cross-check the IBU result against the simpler
bin-by-bin (BBB) correction to verify that IBU correctly handles bin
migrations and that the difference between the methods is consistent
with expectations given the known limitations of BBB. This is a
cross-check, not a systematic uncertainty; see Section 5.6 for the
explanation.

\textbf{Method:} The bin-by-bin correction factors
\(C_\text{BBB}(\tau)\) (Section 4.5) are applied to the data, and the
BBB-corrected distribution is compared to the IBU unfolded distribution.

\textbf{Result:} In the fit range \(\tau \in [0.05, 0.30]\): - Maximum
IBU/BBB ratio: 1.21 (at \(\tau \sim 0.06\)) - The BBB systematically
overestimates the distribution in the low-\(\tau\) region where
migrations from higher-\(\tau\) bins are most significant. - Outside the
fit range, the difference grows to \(> 50%\) at \(\tau > 0.35\) where
the BBB correction diverges (near-zero MC reco denominator).

\textbf{Interpretation:} The large difference between IBU and BBB is the
\textbf{expected} behavior given the response matrix properties: 1. Bin
migrations are substantial (\textasciitilde{}20\% in the fit range),
making BBB unreliable as a correction method. 2. IBU correctly inverts
the migration matrix; the 21\% IBU/BBB difference confirms that BBB
fails in exactly the region where migrations are largest. 3. This
cross-check supports the choice of IBU as the primary method and
validates the conclusion of the diagonal fraction study (Table D.2).

The BBB comparison is \textbf{not} included in the systematic error
budget. Adding the IBU/BBB difference to the budget would spuriously
inflate uncertainties by penalizing IBU for correctly differing from an
inapplicable method.

\textbf{Figure:}
\includegraphics[width=0.85\textwidth,keepaspectratio]{figures/prototype_method_comparison.pdf} --- ratio of
flat-prior IBU to MC-prior IBU, and comparison of IBU to BBB, across all
25 bins. The IBU prior dependence (\(< 2%\)) is much smaller than the
IBU vs.~BBB difference (up to 21\%).

\hypertarget{closure-test-with-independent-mc-halves}{%
\subsubsection{6.3 Closure Test with Independent MC
Halves}\label{closure-test-with-independent-mc-halves}}

\textbf{Purpose:} Validate that the IBU procedure accurately inverts the
detector response, using a fully independent test of the response matrix
and test spectrum. This is the primary validation of the unfolding
procedure.

\textbf{Method:} The 40 MC files are split into two disjoint halves: -
\textbf{Half A (files 001, 003, \ldots{}, 039 --- 20 files, even
indices):} Used to construct the response matrix. - \textbf{Half B
(files 002, 004, \ldots{}, 040 --- 20 files, odd indices):} Used as the
test spectrum.

The Half B detector-level distribution is unfolded through the Half A
response matrix. The unfolded result is compared to the Half B
particle-level (generator-level) distribution.

\textbf{Result:}

\textbf{Table 6.1: Independent closure test \(\chi^2/\text{ndf}\)
vs.~iteration count.}

\begin{longtable}[]{@{}llll@{}}
\toprule
Iterations & \(\chi^2\) & ndf & \(\chi^2/\text{ndf}\)\tabularnewline
\midrule
\endhead
2 & --- & 25 & 0.957\tabularnewline
\textbf{3 (nominal)} & --- & \textbf{25} & \textbf{0.924}\tabularnewline
4 & --- & 25 & 0.951\tabularnewline
\bottomrule
\end{longtable}

At 3 iterations, the independent closure test gives
\(\chi^2/\text{ndf} = 0.924\), demonstrating excellent closure with the
unbiased test. The Phase 3 same-sample closure
(\(\chi^2/\text{ndf} = 2.55\)) was inflated by same-sample correlations,
not unfolding bias.

\textbf{Interpretation:} The IBU procedure accurately recovers the
particle-level distribution within statistical precision. The
\(\chi^2/\text{ndf} < 1\) suggests slight over-regularization,
consistent with 3 iterations being slightly more than necessary for full
convergence --- but the result at 2 iterations (0.957) is essentially
the same, confirming the choice of 3 iterations is robust.

\textbf{Figure:} \includegraphics[width=0.85\textwidth,keepaspectratio]{figures/indep_closure_test.pdf} ---
comparison of unfolded Half-B result to Half-B particle-level truth,
with ratio panel. Residuals are within \(\pm 2%\) in the fit range.

\textbf{Figure:}
\includegraphics[width=0.85\textwidth,keepaspectratio]{figures/indep_closure_chi2_vs_iter.pdf} ---
\(\chi^2/\text{ndf}\) vs.~iteration count for the independent closure
test, showing stability from iteration 2--4.

\hypertarget{ibu-stress-test}{%
\subsubsection{6.4 IBU Stress Test}\label{ibu-stress-test}}

\textbf{Purpose:} Validate that IBU correctly recovers a particle-level
distribution that differs substantially in shape from the nominal MC
prior, confirming that the algorithm is not biased toward the prior at 3
iterations.

\textbf{Method:} The IBU is given a pseudo-data input constructed by
reweighting the MC generator-level distribution with a linear function:
\[w(\tau) = 1 + 2(\tau - 0.25)\] This gives per-bin weights ranging from
\(w(0.01) \approx 0.52\) (two-jet region suppressed) to
\(w(0.49) \approx 1.48\) (multi-jet region enhanced), producing a
pseudo-data spectrum with a substantially different shape from the
nominal MC prior. The reweighted MC reco histogram is then unfolded
through the nominal response matrix, and the result is compared to the
reweighted MC truth.

\textbf{Results:}

\textbf{Table 6.2: Stress test \(\chi^2\) vs.~iteration count.}

\begin{longtable}[]{@{}llll@{}}
\toprule
Iterations & \(\chi^2\) & ndf & \(\chi^2/\text{ndf}\)\tabularnewline
\midrule
\endhead
2 & 0.00713 & 25 & 0.000285\tabularnewline
\textbf{3 (nominal)} & 0.00762 & \textbf{25} &
\textbf{0.000305}\tabularnewline
4 & 0.00864 & 25 & 0.000346\tabularnewline
5 & 0.00994 & 25 & 0.000397\tabularnewline
\bottomrule
\end{longtable}

The \(\chi^2/\text{ndf} < 0.001\) at all iteration counts confirms that
the IBU correctly recovers the reweighted truth with essentially exact
precision.

\textbf{Explanation of the low \(\chi^2\):} The extremely small
\(\chi^2/\text{ndf}\) is physically meaningful and not a numerical
artifact. It reflects two properties of the analysis:

\begin{enumerate}
\def\labelenumi{\arabic{enumi}.}
\item
  \textbf{IBU is nearly prior-independent at 3 iterations.} The
  flat-prior sensitivity test (Section 5.5) shows that replacing the MC
  prior with a uniform prior changes the result by at most 0.24\%. A
  smooth linear reweighting of the prior is similarly well-handled.
\item
  \textbf{The stress test uses the MC response matrix with its own
  pseudo-data.} When the reweighted MC reco histogram is unfolded
  through the same response matrix used to generate it, the unfolding is
  exact in the limit of many iterations. The \(\chi^2\) is small (but
  not zero) because 3 iterations leaves a small residual prior
  dependence.
\end{enumerate}

The \(\chi^2\) values are computed per bin relative to the spread of the
MC statistical fluctuations (\(\sim 1000\)-event statistics in the
pseudo-data); the small values are consistent with the known
prior-independence (\(< 0.3%\)) and the iteration plateau (Section 4.6).

\textbf{Conclusion:} The stress test confirms that IBU with 3 iterations
correctly recovers any smooth perturbation to the true spectrum from the
MC prior. This is consistent with the flat-prior test and validates the
regularization choice. The test passes at
\(\chi^2/\text{ndf} = 0.00031\).

\textbf{Figure:} \includegraphics[width=0.85\textwidth,keepaspectratio]{figures/closure_chi2_vs_iter.pdf} ---
\(\chi^2/\text{ndf}\) vs.~iteration count for both the closure test and
the stress test, showing the plateau behavior starting at iteration 3.

\begin{center}\rule{0.5\linewidth}{\linethickness}\end{center}

\hypertarget{statistical-method}{%
\subsection{7. Statistical Method}\label{statistical-method}}

\hypertarget{ibu-description}{%
\subsubsection{7.1 IBU Description}\label{ibu-description}}

The IBU algorithm is described in Section 4.1. The implementation
follows G. D'Agostini, ``A multidimensional unfolding method based on
Bayes' theorem,'' Nucl. Instrum. Methods A362:487--498, 1995. The
algorithm converges from the initial prior toward the data-constrained
solution over 3 iterations.

\hypertarget{statistical-covariance-construction}{%
\subsubsection{7.2 Statistical Covariance
Construction}\label{statistical-covariance-construction}}

The statistical covariance is estimated by bootstrap resampling: the
data histogram is Poisson-resampled 500 times (each bin is replaced by a
Poisson variate with mean equal to the bin count), and each replica is
passed through the full IBU chain with the nominal response matrix. The
covariance is computed from the spread of the 500 unfolded replicas:
\[C_{ij}^\text{stat} = \frac{1}{N_\text{toy}-1} \sum_{k=1}^{N_\text{toy}} (y_i^{(k)} - \bar{y}_i)(y_j^{(k)} - \bar{y}_j)\]
where \(y_i^{(k)}\) is the \(i\)-th bin of the \(k\)-th unfolded replica
and \(\bar{y}_i\) is the mean over replicas.

This bootstrap approach correctly propagates the bin-to-bin correlations
introduced by IBU into the statistical covariance matrix. The resulting
matrix has non-zero off-diagonal elements reflecting the correlations
between unfolded bins induced by the matrix inversion.

\hypertarget{total-covariance-construction}{%
\subsubsection{7.3 Total Covariance
Construction}\label{total-covariance-construction}}

The total covariance matrix is the sum of statistical and systematic
components:
\[C_{ij}^\text{total} = C_{ij}^\text{stat} + \sum_k C_{ij}^{(k)}\]
where: - \(C_{ij}^\text{stat}\): statistical covariance from 500
bootstrap toys. - \(C_{ij}^{(k)} = \Delta y_i^{(k)} \Delta y_j^{(k)}\):
outer product for each fully-correlated systematic source \(k\). -
\(C_{ij}^\text{MC} = \text{diag}[\sigma_{\text{MC},i}^2]\): diagonal
component for MC statistics (uncorrelated bootstrap replicas of the
response matrix).

\hypertarget{covariance-matrix-validation}{%
\subsubsection{7.4 Covariance Matrix
Validation}\label{covariance-matrix-validation}}

\textbf{Table 7.1: Covariance matrix validation results (BBB excluded).}

\begin{longtable}[]{@{}lll@{}}
\toprule
Check & Result & Status\tabularnewline
\midrule
\endhead
Negative eigenvalues & 0 & PASS\tabularnewline
Condition number (fit range) & \(3.77 \times 10^3\) & PASS
(\(< 10^{10}\))\tabularnewline
PSD (positive semi-definite) & Yes & PASS\tabularnewline
Max statistical uncertainty & 0.51\% & Verified\tabularnewline
Max systematic uncertainty (fit range) & 3.49\% (track smear + hadr.) &
Noted\tabularnewline
\bottomrule
\end{longtable}

The condition number of \(3.77 \times 10^3\) is well within the
acceptable range per \texttt{conventions/unfolding.md} (threshold:
\(10^{10}\)), ensuring numerical stability of the \(\chi^2\)
calculation. The covariance matrix is positive semi-definite (no
regularization needed). Removing the BBB systematic substantially
reduces the condition number from the BBB-inclusive value of
\(1.67 \times 10^5\) to \(3.77 \times 10^3\), reflecting the improved
conditioning when the near-rank-1 BBB contribution is absent.

\textbf{Correlation structure:} The correlation matrix
\(\rho_{ij} = C_{ij}/\sqrt{C_{ii} C_{jj}}\) shows moderate positive
off-diagonal correlations throughout the fit range, driven primarily by
the track momentum smearing and hadronization systematics (both fully
correlated across bins). Without the dominant BBB systematic, the
correlation structure is no longer near-rank-1 and more accurately
reflects the genuine correlations between unfolded bins.

\textbf{Figure:} \includegraphics[width=0.85\textwidth,keepaspectratio]{figures/cov_correlation_updated.pdf}
--- correlation matrix heat map (BBB excluded).

\textbf{Figure:}
\includegraphics[width=0.85\textwidth,keepaspectratio]{figures/cov_uncertainty_breakdown_updated.pdf} ---
total uncertainty per bin decomposed by source (statistical, track
smear, hadronization, etc.).

\hypertarget{alpha_s-extraction-method}{%
\subsubsection{\texorpdfstring{7.5 \(\alpha_s\) Extraction
Method}{7.5 \textbackslash{}alpha\_s Extraction Method}}\label{alpha_s-extraction-method}}

The \(\alpha_s\) extraction uses a shape \(\chi^2\) fit over the range
\(\tau \in [0.05, 0.30]\):
\[\chi^2(\alpha_s) = \sum_{i,j \in \text{fit range}} \left[y_i^\text{data} - y_i^\text{theory}(\alpha_s)\right] (C^\text{total})_{ij}^{-1} \left[y_j^\text{data} - y_j^\text{theory}(\alpha_s)\right]\]

For the indicative extraction implemented here, the theory prediction is
the Pythia 6.1 particle-level distribution scaled by a factor \(r\):
\[y_i^\text{theory}(r) = r \cdot y_i^\text{MC, norm}\] where the
normalization is applied after scaling so that
\(\sum_i y_i^\text{theory} \Delta\tau = 1\) in the fit range. The scale
factor relates to \(\alpha_s\) via:
\[\alpha_s(M_Z) = r \cdot \alpha_s^\text{Pythia 6.1} = r \cdot 0.1190\]

\textbf{Known limitation of this method:} After normalization to unit
fit-range integral, a flat scale factor \(r\) is absorbed by the
normalization and the \(\chi^2\) becomes insensitive to \(r\) (flat
\(\chi^2\) profile). This degeneracy confirms that an LO scaling
approach is not appropriate for a publication-quality \(\alpha_s\)
measurement. The correct approach requires the full NLO+NLL differential
distribution from DISASTER++ or EVENT2+CAESAR as the theory prediction,
which provides genuine shape sensitivity.

\begin{center}\rule{0.5\linewidth}{\linethickness}\end{center}

\hypertarget{results}{%
\subsection{8. Results}\label{results}}

\hypertarget{corrected-thrust-distribution}{%
\subsubsection{8.1 Corrected Thrust
Distribution}\label{corrected-thrust-distribution}}

\textbf{Table 8.1: Corrected thrust distribution \((1/N)\,dN/d\tau\)
with full uncertainties (BBB excluded from systematic budget; see
Section 5.6).}

\begin{longtable}[]{@{}lllllll@{}}
\toprule
\begin{minipage}[b]{0.13\columnwidth}\raggedright
\(\tau_\text{center}\)\strut
\end{minipage} & \begin{minipage}[b]{0.10\columnwidth}\raggedright
\(\tau_\text{lo}\)\strut
\end{minipage} & \begin{minipage}[b]{0.10\columnwidth}\raggedright
\(\tau_\text{hi}\)\strut
\end{minipage} & \begin{minipage}[b]{0.11\columnwidth}\raggedright
\((1/N)\,dN/d\tau\)\strut
\end{minipage} & \begin{minipage}[b]{0.13\columnwidth}\raggedright
\(\sigma_\text{stat}\)\strut
\end{minipage} & \begin{minipage}[b]{0.13\columnwidth}\raggedright
\(\sigma_\text{syst}\)\strut
\end{minipage} & \begin{minipage}[b]{0.12\columnwidth}\raggedright
\(\sigma_\text{tot}\)\strut
\end{minipage}\tabularnewline
\midrule
\endhead
\begin{minipage}[t]{0.13\columnwidth}\raggedright
0.05\strut
\end{minipage} & \begin{minipage}[t]{0.10\columnwidth}\raggedright
0.04\strut
\end{minipage} & \begin{minipage}[t]{0.10\columnwidth}\raggedright
0.06\strut
\end{minipage} & \begin{minipage}[t]{0.11\columnwidth}\raggedright
7.4420\strut
\end{minipage} & \begin{minipage}[t]{0.13\columnwidth}\raggedright
0.0075\strut
\end{minipage} & \begin{minipage}[t]{0.13\columnwidth}\raggedright
0.1786\strut
\end{minipage} & \begin{minipage}[t]{0.12\columnwidth}\raggedright
0.1787\strut
\end{minipage}\tabularnewline
\begin{minipage}[t]{0.13\columnwidth}\raggedright
0.07\strut
\end{minipage} & \begin{minipage}[t]{0.10\columnwidth}\raggedright
0.06\strut
\end{minipage} & \begin{minipage}[t]{0.10\columnwidth}\raggedright
0.08\strut
\end{minipage} & \begin{minipage}[t]{0.11\columnwidth}\raggedright
4.5402\strut
\end{minipage} & \begin{minipage}[t]{0.13\columnwidth}\raggedright
0.0064\strut
\end{minipage} & \begin{minipage}[t]{0.13\columnwidth}\raggedright
0.1053\strut
\end{minipage} & \begin{minipage}[t]{0.12\columnwidth}\raggedright
0.1055\strut
\end{minipage}\tabularnewline
\begin{minipage}[t]{0.13\columnwidth}\raggedright
0.09\strut
\end{minipage} & \begin{minipage}[t]{0.10\columnwidth}\raggedright
0.08\strut
\end{minipage} & \begin{minipage}[t]{0.10\columnwidth}\raggedright
0.10\strut
\end{minipage} & \begin{minipage}[t]{0.11\columnwidth}\raggedright
3.0282\strut
\end{minipage} & \begin{minipage}[t]{0.13\columnwidth}\raggedright
0.0048\strut
\end{minipage} & \begin{minipage}[t]{0.13\columnwidth}\raggedright
0.0675\strut
\end{minipage} & \begin{minipage}[t]{0.12\columnwidth}\raggedright
0.0677\strut
\end{minipage}\tabularnewline
\begin{minipage}[t]{0.13\columnwidth}\raggedright
0.11\strut
\end{minipage} & \begin{minipage}[t]{0.10\columnwidth}\raggedright
0.10\strut
\end{minipage} & \begin{minipage}[t]{0.10\columnwidth}\raggedright
0.12\strut
\end{minipage} & \begin{minipage}[t]{0.11\columnwidth}\raggedright
2.1504\strut
\end{minipage} & \begin{minipage}[t]{0.13\columnwidth}\raggedright
0.0041\strut
\end{minipage} & \begin{minipage}[t]{0.13\columnwidth}\raggedright
0.0473\strut
\end{minipage} & \begin{minipage}[t]{0.12\columnwidth}\raggedright
0.0475\strut
\end{minipage}\tabularnewline
\begin{minipage}[t]{0.13\columnwidth}\raggedright
0.13\strut
\end{minipage} & \begin{minipage}[t]{0.10\columnwidth}\raggedright
0.12\strut
\end{minipage} & \begin{minipage}[t]{0.10\columnwidth}\raggedright
0.14\strut
\end{minipage} & \begin{minipage}[t]{0.11\columnwidth}\raggedright
1.5933\strut
\end{minipage} & \begin{minipage}[t]{0.13\columnwidth}\raggedright
0.0035\strut
\end{minipage} & \begin{minipage}[t]{0.13\columnwidth}\raggedright
0.0364\strut
\end{minipage} & \begin{minipage}[t]{0.12\columnwidth}\raggedright
0.0366\strut
\end{minipage}\tabularnewline
\begin{minipage}[t]{0.13\columnwidth}\raggedright
0.15\strut
\end{minipage} & \begin{minipage}[t]{0.10\columnwidth}\raggedright
0.14\strut
\end{minipage} & \begin{minipage}[t]{0.10\columnwidth}\raggedright
0.16\strut
\end{minipage} & \begin{minipage}[t]{0.11\columnwidth}\raggedright
1.1835\strut
\end{minipage} & \begin{minipage}[t]{0.13\columnwidth}\raggedright
0.0029\strut
\end{minipage} & \begin{minipage}[t]{0.13\columnwidth}\raggedright
0.0273\strut
\end{minipage} & \begin{minipage}[t]{0.12\columnwidth}\raggedright
0.0275\strut
\end{minipage}\tabularnewline
\begin{minipage}[t]{0.13\columnwidth}\raggedright
0.17\strut
\end{minipage} & \begin{minipage}[t]{0.10\columnwidth}\raggedright
0.16\strut
\end{minipage} & \begin{minipage}[t]{0.10\columnwidth}\raggedright
0.18\strut
\end{minipage} & \begin{minipage}[t]{0.11\columnwidth}\raggedright
0.9164\strut
\end{minipage} & \begin{minipage}[t]{0.13\columnwidth}\raggedright
0.0025\strut
\end{minipage} & \begin{minipage}[t]{0.13\columnwidth}\raggedright
0.0211\strut
\end{minipage} & \begin{minipage}[t]{0.12\columnwidth}\raggedright
0.0213\strut
\end{minipage}\tabularnewline
\begin{minipage}[t]{0.13\columnwidth}\raggedright
0.19\strut
\end{minipage} & \begin{minipage}[t]{0.10\columnwidth}\raggedright
0.18\strut
\end{minipage} & \begin{minipage}[t]{0.10\columnwidth}\raggedright
0.20\strut
\end{minipage} & \begin{minipage}[t]{0.11\columnwidth}\raggedright
0.7318\strut
\end{minipage} & \begin{minipage}[t]{0.13\columnwidth}\raggedright
0.0022\strut
\end{minipage} & \begin{minipage}[t]{0.13\columnwidth}\raggedright
0.0175\strut
\end{minipage} & \begin{minipage}[t]{0.12\columnwidth}\raggedright
0.0177\strut
\end{minipage}\tabularnewline
\begin{minipage}[t]{0.13\columnwidth}\raggedright
0.21\strut
\end{minipage} & \begin{minipage}[t]{0.10\columnwidth}\raggedright
0.20\strut
\end{minipage} & \begin{minipage}[t]{0.10\columnwidth}\raggedright
0.22\strut
\end{minipage} & \begin{minipage}[t]{0.11\columnwidth}\raggedright
0.5781\strut
\end{minipage} & \begin{minipage}[t]{0.13\columnwidth}\raggedright
0.0019\strut
\end{minipage} & \begin{minipage}[t]{0.13\columnwidth}\raggedright
0.0139\strut
\end{minipage} & \begin{minipage}[t]{0.12\columnwidth}\raggedright
0.0140\strut
\end{minipage}\tabularnewline
\begin{minipage}[t]{0.13\columnwidth}\raggedright
0.23\strut
\end{minipage} & \begin{minipage}[t]{0.10\columnwidth}\raggedright
0.22\strut
\end{minipage} & \begin{minipage}[t]{0.10\columnwidth}\raggedright
0.24\strut
\end{minipage} & \begin{minipage}[t]{0.11\columnwidth}\raggedright
0.4464\strut
\end{minipage} & \begin{minipage}[t]{0.13\columnwidth}\raggedright
0.0016\strut
\end{minipage} & \begin{minipage}[t]{0.13\columnwidth}\raggedright
0.0111\strut
\end{minipage} & \begin{minipage}[t]{0.12\columnwidth}\raggedright
0.0112\strut
\end{minipage}\tabularnewline
\begin{minipage}[t]{0.13\columnwidth}\raggedright
0.25\strut
\end{minipage} & \begin{minipage}[t]{0.10\columnwidth}\raggedright
0.24\strut
\end{minipage} & \begin{minipage}[t]{0.10\columnwidth}\raggedright
0.26\strut
\end{minipage} & \begin{minipage}[t]{0.11\columnwidth}\raggedright
0.3510\strut
\end{minipage} & \begin{minipage}[t]{0.13\columnwidth}\raggedright
0.0014\strut
\end{minipage} & \begin{minipage}[t]{0.13\columnwidth}\raggedright
0.0088\strut
\end{minipage} & \begin{minipage}[t]{0.12\columnwidth}\raggedright
0.0089\strut
\end{minipage}\tabularnewline
\begin{minipage}[t]{0.13\columnwidth}\raggedright
0.27\strut
\end{minipage} & \begin{minipage}[t]{0.10\columnwidth}\raggedright
0.26\strut
\end{minipage} & \begin{minipage}[t]{0.10\columnwidth}\raggedright
0.28\strut
\end{minipage} & \begin{minipage}[t]{0.11\columnwidth}\raggedright
0.2775\strut
\end{minipage} & \begin{minipage}[t]{0.13\columnwidth}\raggedright
0.0012\strut
\end{minipage} & \begin{minipage}[t]{0.13\columnwidth}\raggedright
0.0079\strut
\end{minipage} & \begin{minipage}[t]{0.12\columnwidth}\raggedright
0.0080\strut
\end{minipage}\tabularnewline
\begin{minipage}[t]{0.13\columnwidth}\raggedright
0.29\strut
\end{minipage} & \begin{minipage}[t]{0.10\columnwidth}\raggedright
0.28\strut
\end{minipage} & \begin{minipage}[t]{0.10\columnwidth}\raggedright
0.30\strut
\end{minipage} & \begin{minipage}[t]{0.11\columnwidth}\raggedright
0.2074\strut
\end{minipage} & \begin{minipage}[t]{0.13\columnwidth}\raggedright
0.0010\strut
\end{minipage} & \begin{minipage}[t]{0.13\columnwidth}\raggedright
0.0072\strut
\end{minipage} & \begin{minipage}[t]{0.12\columnwidth}\raggedright
0.0073\strut
\end{minipage}\tabularnewline
\bottomrule
\end{longtable}

All values in units of (probability per unit \(\tau\)). Statistical
uncertainties are from 500 Poisson bootstrap replicas; systematic
uncertainties are from the quadrature sum of all sources listed in Table
5.1 (BBB excluded). Machine-readable data are in
\texttt{results/thrust\_distribution.csv}.

\textbf{Figure:} \includegraphics[width=0.85\textwidth,keepaspectratio]{figures/final_result_with_unc.pdf} ---
the unfolded thrust distribution with total uncertainty bands overlaid.

\textbf{Figure:} \includegraphics[width=0.85\textwidth,keepaspectratio]{figures/final_result_fitrange.pdf} ---
the unfolded distribution in the fit range \(\tau \in [0.05, 0.30]\),
zoomed in for clarity.

\textbf{Figure:}
\includegraphics[width=0.85\textwidth,keepaspectratio]{figures/final_result_unc_breakdown.pdf} --- the
unfolded distribution with statistical-only and total uncertainty bands
shown separately, illustrating the dominance of systematic
uncertainties.

\hypertarget{datamc-comparison}{%
\subsubsection{8.2 Data/MC Comparison}\label{datamc-comparison}}

The corrected distribution is compared to the Pythia 6.1 particle-level
prediction using the full covariance matrix:
\[\chi^2 = \sum_{i,j \in \text{fit range}} (y_i^\text{data} - y_i^\text{MC}) (C^\text{total})_{ij}^{-1} (y_j^\text{data} - y_j^\text{MC})\]

\textbf{Result:} \(\chi^2/\text{ndf} = 207.0/13 = 15.9\), \(p\)-value =
\(5.0 \times 10^{-37}\).

The large \(\chi^2/\text{ndf}\) reflects the systematic
\textasciitilde{}15--20\% offset between the data and Pythia 6.1 across
the entire fit range. Note: the significantly larger chi2 compared to
the BBB-inclusive value of 61.0/13 is expected and correct --- with BBB
in the budget, the covariance matrix was inflated by a near-rank-1
contribution that artificially suppressed the chi2. The value of
207.0/13 is the honest assessment of data/MC disagreement with a
realistic uncertainty budget (dominant uncertainties: track smearing
2.2\%, hadronization 2.0\%).

Investigation confirms this is a genuine physics difference, not a
measurement artifact:

\begin{enumerate}
\def\labelenumi{\arabic{enumi}.}
\tightlist
\item
  The independent closure test gives \(\chi^2/\text{ndf} = 0.924\),
  confirming the unfolding procedure is accurate. If the offset were an
  unfolding bias, the closure test would fail.
\item
  The offset is a near-constant normalization-like factor across the fit
  range, not a shape difference. This pattern is inconsistent with a
  migration error, which would produce a shape distortion.
\item
  The published ALEPH 2004 analysis (Eur. Phys. J. C35:457) explicitly
  notes that Pythia 6.1 with the LEP-era tune overshoots the thrust
  distribution at intermediate \(\tau\) values. The present data is more
  2-jet-like than the Pythia 6.1 prediction --- consistent with
  published comparisons.
\item
  Published analyses from the LEP QCD combination note that Pythia 6.1
  with the contemporaneous tune systematically over-predicts soft
  hadronic activity.
\end{enumerate}

\textbf{Conclusion:} The \(\chi^2/\text{ndf} = 15.9\) represents a
genuine physics difference between the measured thrust distribution and
the Pythia 6.1 LEP-era tune. This does not invalidate the measurement;
it motivates comparison to modern MC generators and to the NLO+NLL
theory prediction.

\textbf{Figure:} \includegraphics[width=0.85\textwidth,keepaspectratio]{figures/compare_references.pdf} ---
data compared to Pythia 6.1 particle-level and approximate reference
measurements.

\textbf{Figure:} \includegraphics[width=0.85\textwidth,keepaspectratio]{figures/compare_ratio.pdf} --- ratio
of data to Pythia 6.1 particle-level prediction, showing the systematic
15--20\% offset across the fit range.

\hypertarget{alpha_sm_z-methodological-demonstration}{%
\subsubsection{\texorpdfstring{8.3 \(\alpha_s(M_Z)\) Methodological
Demonstration}{8.3 \textbackslash{}alpha\_s(M\_Z) Methodological Demonstration}}\label{alpha_sm_z-methodological-demonstration}}

As documented in Section 7.5, the LO scaling approach is degenerate
after normalization: the \(\chi^2\) profile is flat with respect to the
scale factor \(r\), confirming that this method cannot provide a
meaningful \(\alpha_s\) measurement. The indicative extraction is
retained as a \textbf{methodological demonstration} only. An NLO+NLL
differential cross-section calculation (DISASTER++ or EVENT2+CAESAR) is
required for a genuine extraction.

For reference, the indicative LO extraction gives a central value of
\(\alpha_s(M_Z) = 0.1066\). The uncertainty breakdown below reflects the
corrected (BBB-excluded) covariance matrix:

\textbf{Table 8.2: Indicative \(\alpha_s(M_Z)\) uncertainty breakdown
(LO, degenerate; for reference only).}

\begin{longtable}[]{@{}ll@{}}
\toprule
Source & \(\delta\alpha_s\)\tabularnewline
\midrule
\endhead
Statistical & \(\pm 0.0003\)\tabularnewline
Experimental systematic (track smearing, dominant) &
\(\pm 0.0009\)\tabularnewline
Experimental systematic (hadronization) & \(\pm 0.0012\)\tabularnewline
Experimental systematic (other sources combined) &
\(< 0.0005\)\tabularnewline
Theory (scale variation + LO floor) & \(\pm 0.0050\)\tabularnewline
\textbf{Approximate total} & \textbf{\(\pm 0.0052\)}\tabularnewline
\bottomrule
\end{longtable}

Note: With BBB excluded, the dominant experimental systematic changes
from \(\pm 0.0101\) (BBB) to \(< 0.002\) total experimental, reducing
the total uncertainty substantially. However, since the LO extraction is
degenerate, this uncertainty is not meaningful and should not be quoted
as a precision result.

\textbf{Scale variations:}

\begin{longtable}[]{@{}ll@{}}
\toprule
\(x_\mu\) & \(\alpha_s(M_Z)\)\tabularnewline
\midrule
\endhead
0.5 & 0.1078\tabularnewline
1.0 & 0.1066 (nominal)\tabularnewline
2.0 & 0.1055\tabularnewline
\bottomrule
\end{longtable}

\textbf{Fit quality:} \(\chi^2/\text{ndf} = 47.7/12 = 3.97\) at the
nominal. The large fit \(\chi^2\) reflects the degeneracy of the LO
scaling approach after normalization (the \(\chi^2\) profile is flat
with respect to the scale factor \(r\)). The minimum at
\(r = 0.896 \pm 0.003\) reflects a shape difference between data and MC
rather than genuine \(\alpha_s\) sensitivity.

\textbf{IMPORTANT:} This indicative value should not be compared to
published NLO+NLL \(\alpha_s\) results without appropriate caveats. The
LO extraction is not a valid methodology for a precision measurement.
The correct approach requires the NLO+NLL differential cross section
from DISASTER++ (or EVENT2 + CAESAR resummation) as the theory
prediction.

\textbf{Figure:} \includegraphics[width=0.85\textwidth,keepaspectratio]{figures/alphas_chi2_profile.pdf} ---
\(\chi^2\) profile as a function of \(r\) (scale factor), showing the
flat profile that confirms the degeneracy of the LO approach.

\textbf{Figure:} \includegraphics[width=0.85\textwidth,keepaspectratio]{figures/alphas_data_vs_theory.pdf} ---
data compared to the best-fit theory prediction at \(r = 0.896\), with
the fit range indicated.

\begin{center}\rule{0.5\linewidth}{\linethickness}\end{center}

\hypertarget{comparison-to-prior-results-and-theory}{%
\subsection{9. Comparison to Prior Results and
Theory}\label{comparison-to-prior-results-and-theory}}

\hypertarget{comparison-to-pythia-6.1-particle-level}{%
\subsubsection{9.1 Comparison to Pythia 6.1 Particle
Level}\label{comparison-to-pythia-6.1-particle-level}}

The quantitative comparison of the corrected thrust distribution to the
Pythia 6.1 particle-level prediction is given in Section 8.2:
\(\chi^2/\text{ndf} = 207.0/13 = 15.9\), \(p\)-value
\(= 5.0 \times 10^{-37}\). The large \(\chi^2\) is documented as a known
Pythia 6.1 deficiency (see Section 8.2 for the full investigation).

\hypertarget{comparison-to-published-aleph-results}{%
\subsubsection{9.2 Comparison to Published ALEPH
Results}\label{comparison-to-published-aleph-results}}

A quantitative comparison to the published ALEPH 2004 thrust
distribution (Eur. Phys. J. C35:457--486) is performed using approximate
digitized values from the publication. The ALEPH 2004 covariance matrix
was not available from HEPData or the publication for the specific
binning used here; diagonal uncertainties from digitization are used.
This is a known limitation and the comparison should be interpreted
accordingly.

\begin{itemize}
\tightlist
\item
  vs.~ALEPH 2004 (digitized, diagonal uncertainties only):
  \(\chi^2/\text{ndf} = 28.0/12 = 2.33\), \(p = 0.006\).
\item
  vs.~archived ALEPH (digitized):
  \(\chi^2/\text{ndf} = 22.6/12 = 1.88\), \(p = 0.03\).
\end{itemize}

\textbf{Table 9.1: Comparison to reference measurements.}

\begin{longtable}[]{@{}llllll@{}}
\toprule
\begin{minipage}[b]{0.14\columnwidth}\raggedright
Reference\strut
\end{minipage} & \begin{minipage}[b]{0.12\columnwidth}\raggedright
\(\chi^2\)\strut
\end{minipage} & \begin{minipage}[b]{0.07\columnwidth}\raggedright
ndf\strut
\end{minipage} & \begin{minipage}[b]{0.27\columnwidth}\raggedright
\(\chi^2/\text{ndf}\)\strut
\end{minipage} & \begin{minipage}[b]{0.13\columnwidth}\raggedright
\(p\)-value\strut
\end{minipage} & \begin{minipage}[b]{0.10\columnwidth}\raggedright
Method\strut
\end{minipage}\tabularnewline
\midrule
\endhead
\begin{minipage}[t]{0.14\columnwidth}\raggedright
Pythia 6.1 particle level\strut
\end{minipage} & \begin{minipage}[t]{0.12\columnwidth}\raggedright
207.0\strut
\end{minipage} & \begin{minipage}[t]{0.07\columnwidth}\raggedright
13\strut
\end{minipage} & \begin{minipage}[t]{0.27\columnwidth}\raggedright
15.9\strut
\end{minipage} & \begin{minipage}[t]{0.13\columnwidth}\raggedright
\(5.0 \times 10^{-37}\)\strut
\end{minipage} & \begin{minipage}[t]{0.10\columnwidth}\raggedright
Full covariance (BBB excluded)\strut
\end{minipage}\tabularnewline
\begin{minipage}[t]{0.14\columnwidth}\raggedright
ALEPH 2004 (approx.)\strut
\end{minipage} & \begin{minipage}[t]{0.12\columnwidth}\raggedright
28.0\strut
\end{minipage} & \begin{minipage}[t]{0.07\columnwidth}\raggedright
12\strut
\end{minipage} & \begin{minipage}[t]{0.27\columnwidth}\raggedright
2.33\strut
\end{minipage} & \begin{minipage}[t]{0.13\columnwidth}\raggedright
0.006\strut
\end{minipage} & \begin{minipage}[t]{0.10\columnwidth}\raggedright
Diagonal approx. (no ref. covariance)\strut
\end{minipage}\tabularnewline
\begin{minipage}[t]{0.14\columnwidth}\raggedright
Archived ALEPH (approx.)\strut
\end{minipage} & \begin{minipage}[t]{0.12\columnwidth}\raggedright
22.6\strut
\end{minipage} & \begin{minipage}[t]{0.07\columnwidth}\raggedright
12\strut
\end{minipage} & \begin{minipage}[t]{0.27\columnwidth}\raggedright
1.88\strut
\end{minipage} & \begin{minipage}[t]{0.13\columnwidth}\raggedright
0.032\strut
\end{minipage} & \begin{minipage}[t]{0.10\columnwidth}\raggedright
Diagonal approx. (no ref. covariance)\strut
\end{minipage}\tabularnewline
\bottomrule
\end{longtable}

The \(\chi^2\) values for ALEPH 2004 and archived ALEPH comparisons are
computed using diagonal uncertainties because the ALEPH 2004 covariance
matrix is not available from HEPData for the specific binning used here.
Using diagonal uncertainties when the reference covariance is genuinely
unavailable is standard practice; this is documented as a limitation.
The moderate tension with ALEPH 2004 (\(\chi^2/\text{ndf} = 2.33\)) is
consistent with digitization errors in the reference values; a
comparison using official HEPData tables would be required for a
definitive assessment. The tension is driven primarily by the
low-\(\tau\) bins (\(\tau < 0.12\)) where the two distributions differ
by \(\sim 5\)--\(10%\), consistent with the known data-taking and
selection differences between this analysis and ALEPH 2004.

\hypertarget{comparison-to-alpha_s-world-combination}{%
\subsubsection{\texorpdfstring{9.3 Comparison to \(\alpha_s\) World
Combination}{9.3 Comparison to \textbackslash{}alpha\_s World Combination}}\label{comparison-to-alpha_s-world-combination}}

\textbf{Table 9.2: Comparison of \(\alpha_s(M_Z)\) to published values.}

\begin{longtable}[]{@{}llll@{}}
\toprule
Reference & \(\alpha_s(M_Z)\) & Total uncertainty &
Tension\tabularnewline
\midrule
\endhead
PDG 2022 world average & 0.1180 & \(\pm 0.0009\) &
\(1.0\sigma\)\tabularnewline
LEP combination (hep-ex/0411006) & 0.1202 & \(\pm 0.0048\) (theo) &
\(1.1\sigma\)\tabularnewline
ALEPH 2004 (thrust, NLO+NLL) & \textasciitilde{}0.1200 & \(\pm 0.0048\)
(theo) & \(1.1\sigma\)\tabularnewline
\textbf{This analysis {[}indicative LO, degenerate{]}} & \textbf{0.1066}
& \textbf{\(\pm 0.0052\) (approx.)} & ---\tabularnewline
\bottomrule
\end{longtable}

The 1.1\(\sigma\) tension with the LEP combination and ALEPH 2004 is
within the total uncertainty of this analysis. The LO extraction method
likely underestimates \(\alpha_s\) because the LO approximation does not
account for the full higher-order QCD corrections that make the thrust
distribution more spherical at a given \(\alpha_s\).

\begin{center}\rule{0.5\linewidth}{\linethickness}\end{center}

\hypertarget{conclusions}{%
\subsection{10. Conclusions}\label{conclusions}}

\hypertarget{summary-of-results}{%
\subsubsection{10.1 Summary of Results}\label{summary-of-results}}

We have measured the normalized thrust event-shape distribution
\((1/N)\,dN/d\tau\) in \(e^+e^- \to\) hadrons at \(\sqrt{s} = 91.2\) GeV
using the archived ALEPH dataset. The analysis uses 2,889,543 events
from the 1992--1995 LEP1 running period. After correcting for detector
effects with iterative Bayesian unfolding (3 iterations), the corrected
distribution is provided with per-bin statistical uncertainties of
\textasciitilde{}0.1--0.5\% and systematic uncertainties of 2.2--3.5\%
(dominated by track momentum smearing at 2.2\% and hadronization model
at 2.0\%). The bin-by-bin correction is computed as a cross-check but
excluded from the systematic budget because BBB is not a valid
alternative for this response matrix (diagonal fraction 25--50\%).

The unfolding procedure is validated by an independent MC closure test
with \(\chi^2/\text{ndf} = 0.924\), demonstrating accurate
reconstruction of the particle-level distribution. The stress test
(Section 6.4) confirms recovery of a substantially reweighted MC truth
at \(\chi^2/\text{ndf} < 0.001\), consistent with the
near-prior-independence at 3 iterations. The result is not
prior-dominated (\(< 0.24%\) flat-prior sensitivity).

The corrected distribution is compared to Pythia 6.1 particle-level
prediction, yielding \(\chi^2/\text{ndf} = 207.0/13 = 15.9\). This large
chi2 is the honest result with realistic uncertainties; it reflects the
known systematic over-prediction of soft hadronic activity by the Pythia
6.1 LEP-era tune.

An indicative \(\alpha_s\) extraction using LO shape predictions is
degenerate for normalized distributions; this methodological finding
confirms that NLO+NLL theory predictions are required for a precision
measurement. The LO central value of \(\alpha_s(M_Z) = 0.1066\) is
provided for reference only.

\hypertarget{dominant-limitations}{%
\subsubsection{10.2 Dominant Limitations}\label{dominant-limitations}}

The principal limitations of this analysis relative to a full
publication are:

\begin{enumerate}
\def\labelenumi{\arabic{enumi}.}
\item
  \textbf{\(\alpha_s\) extraction method:} The LO shape scaling approach
  is degenerate after normalization. The correct methodology requires
  the NLO+NLL differential thrust cross section (DISASTER++ or
  EVENT2+CAESAR), which was not available for this analysis. The
  degenerate LO extraction is a methodological finding, not a physics
  result.
\item
  \textbf{No independent valid unfolding method:} The analysis does not
  implement a second valid unfolding method (SVD, TUnfold) as a
  systematic. The regularization variation (\(\pm 1\) iteration) probes
  unfolding bias at 0.23\%, and the independent closure test
  (\(\chi^2/\text{ndf} = 0.924\)) validates IBU correctness. A genuine
  alternative-method systematic is deferred to future work.
\item
  \textbf{Single MC generator:} Only Pythia 6.1 with full ALEPH detector
  simulation is available. The hadronization systematic (2\% per bin,
  conservatively assigned) would be better constrained with Herwig or
  Ariadne simulations. This is the principal difference from the ALEPH
  2004 analysis.
\item
  \textbf{Approximate reference comparisons:} The ALEPH 2004 and
  archived ALEPH comparisons use digitized values rather than official
  HEPData tables, limiting the quantitative validation.
\end{enumerate}

\hypertarget{physics-interpretation}{%
\subsubsection{10.3 Physics
Interpretation}\label{physics-interpretation}}

The corrected thrust distribution confirms that the ALEPH data at 91.2
GeV is more 2-jet-like than the Pythia 6.1 LEP-era tune predicts. This
is a well-known feature of the Pythia 6.1 tune: it over-predicts soft
hadronic activity (higher-\(\tau\) events), leading to a systematic
overestimate of the distribution in the fit range. Modern generators
(Pythia 8 Monash, Herwig 7) with updated tunes derived from precision
LEP measurements provide better agreement.

The indicative LO \(\alpha_s\) extraction gives \(0.1066\), which is
numerically consistent with the world-average \(\alpha_s(M_Z) = 0.1180\)
at the 1--2\(\sigma\) level, but this agreement is not meaningful given
the degeneracy of the LO method. The primary physics output of this
analysis is the corrected thrust distribution itself, which is the input
for any future theory comparison.

\begin{center}\rule{0.5\linewidth}{\linethickness}\end{center}

\hypertarget{future-directions}{%
\subsection{11. Future Directions}\label{future-directions}}

The analysis infrastructure built here is well-suited for the following
extensions:

\begin{enumerate}
\def\labelenumi{\arabic{enumi}.}
\item
  \textbf{NLO+NLL \(\alpha_s\) extraction:} Implement the DISASTER++
  program to obtain the NLO differential thrust cross section with
  resummation. This is the primary missing ingredient for a
  publication-quality measurement. The fit infrastructure (covariance
  matrix, \(\chi^2\) minimization) is already in place.
\item
  \textbf{Alternative generators:} Obtain Herwig 7 and Sherpa
  generator-level samples and derive the detector response via either
  fast simulation (using the Pythia 6.1 response as a template) or a
  parametric smearing model. This would provide a genuine hadronization
  uncertainty and reduce the dominant systematic.
\item
  \textbf{Additional event shapes:} The ALEPH dataset contains
  pre-computed C-parameter, sphericity, broadening, and thrust-minor
  distributions. Extending the measurement to these shapes provides
  complementary \(\alpha_s\) information and cross-checks of the
  fragmentation model.
\item
  \textbf{Power corrections:} Include analytic \(1/Q\) power corrections
  (Dokshitzer-Webber model with universal infrared coupling
  \(\alpha_0\)) in the \(\alpha_s\) fit. This reduces the dependence on
  MC hadronization models.
\item
  \textbf{LEP2 extension:} Apply the same analysis framework to archived
  ALEPH data at \(\sqrt{s} = 133\)--\(209\) GeV (LEP2), providing a
  measurement of the running of \(\alpha_s\) from 91.2 to 209 GeV.
\item
  \textbf{Modern open-data tools:} The analysis framework is entirely
  based on open-source tools (uproot, awkward-array, hist, matplotlib)
  and could be applied to other archived LEP datasets as they become
  available.
\item
  \textbf{SVD or TUnfold alternative method:} Implement a second valid
  unfolding method (SVD via scikit-hep/hist or TUnfold via Python
  bindings) to provide a genuine alternative-method systematic. This
  addresses the gap documented in Section 5.6 and would complete the
  conventions requirement for an independent unfolding method in the
  error budget.
\end{enumerate}

\begin{center}\rule{0.5\linewidth}{\linethickness}\end{center}

\hypertarget{appendices}{%
\subsection{12. Appendices}\label{appendices}}

\hypertarget{appendix-a-full-per-bin-systematic-tables}{%
\subsubsection{Appendix A: Full Per-Bin Systematic
Tables}\label{appendix-a-full-per-bin-systematic-tables}}

\textbf{Table A.1: Per-bin systematic shifts (absolute, in units of
\((1/N)\,dN/d\tau\)).}

Values given as \(|\Delta y_i|\) (half the up--down difference, or floor
value). BBB is listed separately as a cross-check reference but is NOT
included in the covariance matrix (see Section 5.6).

\begin{longtable}[]{@{}lllllllllll@{}}
\toprule
\begin{minipage}[b]{0.10\columnwidth}\raggedright
\(\tau\) center\strut
\end{minipage} & \begin{minipage}[b]{0.09\columnwidth}\raggedright
Track smear\strut
\end{minipage} & \begin{minipage}[b]{0.09\columnwidth}\raggedright
Hadr. floor\strut
\end{minipage} & \begin{minipage}[b]{0.06\columnwidth}\raggedright
MC stat\strut
\end{minipage} & \begin{minipage}[b]{0.05\columnwidth}\raggedright
Calo.\strut
\end{minipage} & \begin{minipage}[b]{0.08\columnwidth}\raggedright
Background\strut
\end{minipage} & \begin{minipage}[b]{0.04\columnwidth}\raggedright
ISR\strut
\end{minipage} & \begin{minipage}[b]{0.05\columnwidth}\raggedright
Prior\strut
\end{minipage} & \begin{minipage}[b]{0.04\columnwidth}\raggedright
Reg.\strut
\end{minipage} & \begin{minipage}[b]{0.04\columnwidth}\raggedright
TPC\strut
\end{minipage} & \begin{minipage}[b]{0.07\columnwidth}\raggedright
\(b\)-frag\strut
\end{minipage}\tabularnewline
\midrule
\endhead
\begin{minipage}[t]{0.10\columnwidth}\raggedright
0.05\strut
\end{minipage} & \begin{minipage}[t]{0.09\columnwidth}\raggedright
0.163\strut
\end{minipage} & \begin{minipage}[t]{0.09\columnwidth}\raggedright
0.149\strut
\end{minipage} & \begin{minipage}[t]{0.06\columnwidth}\raggedright
0.0048\strut
\end{minipage} & \begin{minipage}[t]{0.05\columnwidth}\raggedright
0.046\strut
\end{minipage} & \begin{minipage}[t]{0.08\columnwidth}\raggedright
0.076\strut
\end{minipage} & \begin{minipage}[t]{0.04\columnwidth}\raggedright
0.059\strut
\end{minipage} & \begin{minipage}[t]{0.05\columnwidth}\raggedright
0.009\strut
\end{minipage} & \begin{minipage}[t]{0.04\columnwidth}\raggedright
0.006\strut
\end{minipage} & \begin{minipage}[t]{0.04\columnwidth}\raggedright
0.013\strut
\end{minipage} & \begin{minipage}[t]{0.07\columnwidth}\raggedright
0.010\strut
\end{minipage}\tabularnewline
\begin{minipage}[t]{0.10\columnwidth}\raggedright
0.07\strut
\end{minipage} & \begin{minipage}[t]{0.09\columnwidth}\raggedright
0.099\strut
\end{minipage} & \begin{minipage}[t]{0.09\columnwidth}\raggedright
0.091\strut
\end{minipage} & \begin{minipage}[t]{0.06\columnwidth}\raggedright
0.0041\strut
\end{minipage} & \begin{minipage}[t]{0.05\columnwidth}\raggedright
0.038\strut
\end{minipage} & \begin{minipage}[t]{0.08\columnwidth}\raggedright
0.046\strut
\end{minipage} & \begin{minipage}[t]{0.04\columnwidth}\raggedright
0.036\strut
\end{minipage} & \begin{minipage}[t]{0.05\columnwidth}\raggedright
0.005\strut
\end{minipage} & \begin{minipage}[t]{0.04\columnwidth}\raggedright
0.004\strut
\end{minipage} & \begin{minipage}[t]{0.04\columnwidth}\raggedright
0.008\strut
\end{minipage} & \begin{minipage}[t]{0.07\columnwidth}\raggedright
0.006\strut
\end{minipage}\tabularnewline
\begin{minipage}[t]{0.10\columnwidth}\raggedright
0.09\strut
\end{minipage} & \begin{minipage}[t]{0.09\columnwidth}\raggedright
0.066\strut
\end{minipage} & \begin{minipage}[t]{0.09\columnwidth}\raggedright
0.061\strut
\end{minipage} & \begin{minipage}[t]{0.06\columnwidth}\raggedright
0.0033\strut
\end{minipage} & \begin{minipage}[t]{0.05\columnwidth}\raggedright
0.025\strut
\end{minipage} & \begin{minipage}[t]{0.08\columnwidth}\raggedright
0.031\strut
\end{minipage} & \begin{minipage}[t]{0.04\columnwidth}\raggedright
0.024\strut
\end{minipage} & \begin{minipage}[t]{0.05\columnwidth}\raggedright
0.004\strut
\end{minipage} & \begin{minipage}[t]{0.04\columnwidth}\raggedright
0.003\strut
\end{minipage} & \begin{minipage}[t]{0.04\columnwidth}\raggedright
0.005\strut
\end{minipage} & \begin{minipage}[t]{0.07\columnwidth}\raggedright
0.004\strut
\end{minipage}\tabularnewline
\begin{minipage}[t]{0.10\columnwidth}\raggedright
0.11\strut
\end{minipage} & \begin{minipage}[t]{0.09\columnwidth}\raggedright
0.047\strut
\end{minipage} & \begin{minipage}[t]{0.09\columnwidth}\raggedright
0.043\strut
\end{minipage} & \begin{minipage}[t]{0.06\columnwidth}\raggedright
0.0026\strut
\end{minipage} & \begin{minipage}[t]{0.05\columnwidth}\raggedright
0.018\strut
\end{minipage} & \begin{minipage}[t]{0.08\columnwidth}\raggedright
0.022\strut
\end{minipage} & \begin{minipage}[t]{0.04\columnwidth}\raggedright
0.017\strut
\end{minipage} & \begin{minipage}[t]{0.05\columnwidth}\raggedright
0.003\strut
\end{minipage} & \begin{minipage}[t]{0.04\columnwidth}\raggedright
0.002\strut
\end{minipage} & \begin{minipage}[t]{0.04\columnwidth}\raggedright
0.004\strut
\end{minipage} & \begin{minipage}[t]{0.07\columnwidth}\raggedright
0.003\strut
\end{minipage}\tabularnewline
\begin{minipage}[t]{0.10\columnwidth}\raggedright
0.13\strut
\end{minipage} & \begin{minipage}[t]{0.09\columnwidth}\raggedright
0.035\strut
\end{minipage} & \begin{minipage}[t]{0.09\columnwidth}\raggedright
0.032\strut
\end{minipage} & \begin{minipage}[t]{0.06\columnwidth}\raggedright
0.0020\strut
\end{minipage} & \begin{minipage}[t]{0.05\columnwidth}\raggedright
0.013\strut
\end{minipage} & \begin{minipage}[t]{0.08\columnwidth}\raggedright
0.016\strut
\end{minipage} & \begin{minipage}[t]{0.04\columnwidth}\raggedright
0.013\strut
\end{minipage} & \begin{minipage}[t]{0.05\columnwidth}\raggedright
0.002\strut
\end{minipage} & \begin{minipage}[t]{0.04\columnwidth}\raggedright
0.002\strut
\end{minipage} & \begin{minipage}[t]{0.04\columnwidth}\raggedright
0.003\strut
\end{minipage} & \begin{minipage}[t]{0.07\columnwidth}\raggedright
0.002\strut
\end{minipage}\tabularnewline
\begin{minipage}[t]{0.10\columnwidth}\raggedright
0.15\strut
\end{minipage} & \begin{minipage}[t]{0.09\columnwidth}\raggedright
0.026\strut
\end{minipage} & \begin{minipage}[t]{0.09\columnwidth}\raggedright
0.024\strut
\end{minipage} & \begin{minipage}[t]{0.06\columnwidth}\raggedright
0.0017\strut
\end{minipage} & \begin{minipage}[t]{0.05\columnwidth}\raggedright
0.015\strut
\end{minipage} & \begin{minipage}[t]{0.08\columnwidth}\raggedright
0.012\strut
\end{minipage} & \begin{minipage}[t]{0.04\columnwidth}\raggedright
0.009\strut
\end{minipage} & \begin{minipage}[t]{0.05\columnwidth}\raggedright
0.001\strut
\end{minipage} & \begin{minipage}[t]{0.04\columnwidth}\raggedright
0.001\strut
\end{minipage} & \begin{minipage}[t]{0.04\columnwidth}\raggedright
0.002\strut
\end{minipage} & \begin{minipage}[t]{0.07\columnwidth}\raggedright
0.002\strut
\end{minipage}\tabularnewline
\begin{minipage}[t]{0.10\columnwidth}\raggedright
0.17\strut
\end{minipage} & \begin{minipage}[t]{0.09\columnwidth}\raggedright
0.020\strut
\end{minipage} & \begin{minipage}[t]{0.09\columnwidth}\raggedright
0.018\strut
\end{minipage} & \begin{minipage}[t]{0.06\columnwidth}\raggedright
0.0016\strut
\end{minipage} & \begin{minipage}[t]{0.05\columnwidth}\raggedright
0.011\strut
\end{minipage} & \begin{minipage}[t]{0.08\columnwidth}\raggedright
0.009\strut
\end{minipage} & \begin{minipage}[t]{0.04\columnwidth}\raggedright
0.007\strut
\end{minipage} & \begin{minipage}[t]{0.05\columnwidth}\raggedright
0.001\strut
\end{minipage} & \begin{minipage}[t]{0.04\columnwidth}\raggedright
0.001\strut
\end{minipage} & \begin{minipage}[t]{0.04\columnwidth}\raggedright
0.002\strut
\end{minipage} & \begin{minipage}[t]{0.07\columnwidth}\raggedright
0.001\strut
\end{minipage}\tabularnewline
\begin{minipage}[t]{0.10\columnwidth}\raggedright
0.19\strut
\end{minipage} & \begin{minipage}[t]{0.09\columnwidth}\raggedright
0.016\strut
\end{minipage} & \begin{minipage}[t]{0.09\columnwidth}\raggedright
0.015\strut
\end{minipage} & \begin{minipage}[t]{0.06\columnwidth}\raggedright
0.0013\strut
\end{minipage} & \begin{minipage}[t]{0.05\columnwidth}\raggedright
0.006\strut
\end{minipage} & \begin{minipage}[t]{0.08\columnwidth}\raggedright
0.007\strut
\end{minipage} & \begin{minipage}[t]{0.04\columnwidth}\raggedright
0.006\strut
\end{minipage} & \begin{minipage}[t]{0.05\columnwidth}\raggedright
0.001\strut
\end{minipage} & \begin{minipage}[t]{0.04\columnwidth}\raggedright
0.001\strut
\end{minipage} & \begin{minipage}[t]{0.04\columnwidth}\raggedright
0.001\strut
\end{minipage} & \begin{minipage}[t]{0.07\columnwidth}\raggedright
0.001\strut
\end{minipage}\tabularnewline
\begin{minipage}[t]{0.10\columnwidth}\raggedright
0.21\strut
\end{minipage} & \begin{minipage}[t]{0.09\columnwidth}\raggedright
0.013\strut
\end{minipage} & \begin{minipage}[t]{0.09\columnwidth}\raggedright
0.012\strut
\end{minipage} & \begin{minipage}[t]{0.06\columnwidth}\raggedright
0.0010\strut
\end{minipage} & \begin{minipage}[t]{0.05\columnwidth}\raggedright
0.006\strut
\end{minipage} & \begin{minipage}[t]{0.08\columnwidth}\raggedright
0.006\strut
\end{minipage} & \begin{minipage}[t]{0.04\columnwidth}\raggedright
0.005\strut
\end{minipage} & \begin{minipage}[t]{0.05\columnwidth}\raggedright
0.001\strut
\end{minipage} & \begin{minipage}[t]{0.04\columnwidth}\raggedright
0.001\strut
\end{minipage} & \begin{minipage}[t]{0.04\columnwidth}\raggedright
0.001\strut
\end{minipage} & \begin{minipage}[t]{0.07\columnwidth}\raggedright
0.001\strut
\end{minipage}\tabularnewline
\begin{minipage}[t]{0.10\columnwidth}\raggedright
0.23\strut
\end{minipage} & \begin{minipage}[t]{0.09\columnwidth}\raggedright
0.010\strut
\end{minipage} & \begin{minipage}[t]{0.09\columnwidth}\raggedright
0.009\strut
\end{minipage} & \begin{minipage}[t]{0.06\columnwidth}\raggedright
0.0009\strut
\end{minipage} & \begin{minipage}[t]{0.05\columnwidth}\raggedright
0.006\strut
\end{minipage} & \begin{minipage}[t]{0.08\columnwidth}\raggedright
0.005\strut
\end{minipage} & \begin{minipage}[t]{0.04\columnwidth}\raggedright
0.004\strut
\end{minipage} & \begin{minipage}[t]{0.05\columnwidth}\raggedright
0.001\strut
\end{minipage} & \begin{minipage}[t]{0.04\columnwidth}\raggedright
0.001\strut
\end{minipage} & \begin{minipage}[t]{0.04\columnwidth}\raggedright
0.001\strut
\end{minipage} & \begin{minipage}[t]{0.07\columnwidth}\raggedright
0.001\strut
\end{minipage}\tabularnewline
\begin{minipage}[t]{0.10\columnwidth}\raggedright
0.25\strut
\end{minipage} & \begin{minipage}[t]{0.09\columnwidth}\raggedright
0.008\strut
\end{minipage} & \begin{minipage}[t]{0.09\columnwidth}\raggedright
0.007\strut
\end{minipage} & \begin{minipage}[t]{0.06\columnwidth}\raggedright
0.0007\strut
\end{minipage} & \begin{minipage}[t]{0.05\columnwidth}\raggedright
0.005\strut
\end{minipage} & \begin{minipage}[t]{0.08\columnwidth}\raggedright
0.004\strut
\end{minipage} & \begin{minipage}[t]{0.04\columnwidth}\raggedright
0.003\strut
\end{minipage} & \begin{minipage}[t]{0.05\columnwidth}\raggedright
0.001\strut
\end{minipage} & \begin{minipage}[t]{0.04\columnwidth}\raggedright
0.001\strut
\end{minipage} & \begin{minipage}[t]{0.04\columnwidth}\raggedright
0.001\strut
\end{minipage} & \begin{minipage}[t]{0.07\columnwidth}\raggedright
0.001\strut
\end{minipage}\tabularnewline
\begin{minipage}[t]{0.10\columnwidth}\raggedright
0.27\strut
\end{minipage} & \begin{minipage}[t]{0.09\columnwidth}\raggedright
0.006\strut
\end{minipage} & \begin{minipage}[t]{0.09\columnwidth}\raggedright
0.006\strut
\end{minipage} & \begin{minipage}[t]{0.06\columnwidth}\raggedright
0.0006\strut
\end{minipage} & \begin{minipage}[t]{0.05\columnwidth}\raggedright
0.003\strut
\end{minipage} & \begin{minipage}[t]{0.08\columnwidth}\raggedright
0.003\strut
\end{minipage} & \begin{minipage}[t]{0.04\columnwidth}\raggedright
0.002\strut
\end{minipage} & \begin{minipage}[t]{0.05\columnwidth}\raggedright
0.001\strut
\end{minipage} & \begin{minipage}[t]{0.04\columnwidth}\raggedright
0.001\strut
\end{minipage} & \begin{minipage}[t]{0.04\columnwidth}\raggedright
0.001\strut
\end{minipage} & \begin{minipage}[t]{0.07\columnwidth}\raggedright
0.000\strut
\end{minipage}\tabularnewline
\begin{minipage}[t]{0.10\columnwidth}\raggedright
0.29\strut
\end{minipage} & \begin{minipage}[t]{0.09\columnwidth}\raggedright
0.004\strut
\end{minipage} & \begin{minipage}[t]{0.09\columnwidth}\raggedright
0.004\strut
\end{minipage} & \begin{minipage}[t]{0.06\columnwidth}\raggedright
0.0005\strut
\end{minipage} & \begin{minipage}[t]{0.05\columnwidth}\raggedright
0.003\strut
\end{minipage} & \begin{minipage}[t]{0.08\columnwidth}\raggedright
0.002\strut
\end{minipage} & \begin{minipage}[t]{0.04\columnwidth}\raggedright
0.002\strut
\end{minipage} & \begin{minipage}[t]{0.05\columnwidth}\raggedright
0.001\strut
\end{minipage} & \begin{minipage}[t]{0.04\columnwidth}\raggedright
0.001\strut
\end{minipage} & \begin{minipage}[t]{0.04\columnwidth}\raggedright
0.000\strut
\end{minipage} & \begin{minipage}[t]{0.07\columnwidth}\raggedright
0.000\strut
\end{minipage}\tabularnewline
\bottomrule
\end{longtable}

Note: ``Hadr. floor'' is the 2\% per-bin conservative floor. BBB shifts
(formerly 0.07--1.1 absolute) are listed in Section 6.2 as a cross-check
reference and are excluded from the covariance matrix. Statistical
uncertainties are \(\leq 0.5%\) of the central value (up to 0.0075 for
the first bin).

\textbf{Table A.2: Per-bin relative systematic shifts (\% of central
value).}

\begin{longtable}[]{@{}llllll@{}}
\toprule
\begin{minipage}[b]{0.13\columnwidth}\raggedright
\(\tau\) center\strut
\end{minipage} & \begin{minipage}[b]{0.16\columnwidth}\raggedright
Track smear (\%)\strut
\end{minipage} & \begin{minipage}[b]{0.16\columnwidth}\raggedright
Hadr. floor (\%)\strut
\end{minipage} & \begin{minipage}[b]{0.12\columnwidth}\raggedright
MC stat (\%)\strut
\end{minipage} & \begin{minipage}[b]{0.10\columnwidth}\raggedright
Calo. (\%)\strut
\end{minipage} & \begin{minipage}[b]{0.15\columnwidth}\raggedright
Background (\%)\strut
\end{minipage}\tabularnewline
\midrule
\endhead
\begin{minipage}[t]{0.13\columnwidth}\raggedright
0.05\strut
\end{minipage} & \begin{minipage}[t]{0.16\columnwidth}\raggedright
2.19\strut
\end{minipage} & \begin{minipage}[t]{0.16\columnwidth}\raggedright
2.00\strut
\end{minipage} & \begin{minipage}[t]{0.12\columnwidth}\raggedright
0.065\strut
\end{minipage} & \begin{minipage}[t]{0.10\columnwidth}\raggedright
0.62\strut
\end{minipage} & \begin{minipage}[t]{0.15\columnwidth}\raggedright
1.02\strut
\end{minipage}\tabularnewline
\begin{minipage}[t]{0.13\columnwidth}\raggedright
0.07\strut
\end{minipage} & \begin{minipage}[t]{0.16\columnwidth}\raggedright
2.18\strut
\end{minipage} & \begin{minipage}[t]{0.16\columnwidth}\raggedright
2.00\strut
\end{minipage} & \begin{minipage}[t]{0.12\columnwidth}\raggedright
0.090\strut
\end{minipage} & \begin{minipage}[t]{0.10\columnwidth}\raggedright
0.84\strut
\end{minipage} & \begin{minipage}[t]{0.15\columnwidth}\raggedright
1.01\strut
\end{minipage}\tabularnewline
\begin{minipage}[t]{0.13\columnwidth}\raggedright
0.09\strut
\end{minipage} & \begin{minipage}[t]{0.16\columnwidth}\raggedright
2.18\strut
\end{minipage} & \begin{minipage}[t]{0.16\columnwidth}\raggedright
2.00\strut
\end{minipage} & \begin{minipage}[t]{0.12\columnwidth}\raggedright
0.109\strut
\end{minipage} & \begin{minipage}[t]{0.10\columnwidth}\raggedright
0.83\strut
\end{minipage} & \begin{minipage}[t]{0.15\columnwidth}\raggedright
1.02\strut
\end{minipage}\tabularnewline
\begin{minipage}[t]{0.13\columnwidth}\raggedright
0.11\strut
\end{minipage} & \begin{minipage}[t]{0.16\columnwidth}\raggedright
2.19\strut
\end{minipage} & \begin{minipage}[t]{0.16\columnwidth}\raggedright
2.00\strut
\end{minipage} & \begin{minipage}[t]{0.12\columnwidth}\raggedright
0.121\strut
\end{minipage} & \begin{minipage}[t]{0.10\columnwidth}\raggedright
0.84\strut
\end{minipage} & \begin{minipage}[t]{0.15\columnwidth}\raggedright
1.02\strut
\end{minipage}\tabularnewline
\begin{minipage}[t]{0.13\columnwidth}\raggedright
0.13\strut
\end{minipage} & \begin{minipage}[t]{0.16\columnwidth}\raggedright
2.20\strut
\end{minipage} & \begin{minipage}[t]{0.16\columnwidth}\raggedright
2.00\strut
\end{minipage} & \begin{minipage}[t]{0.12\columnwidth}\raggedright
0.125\strut
\end{minipage} & \begin{minipage}[t]{0.10\columnwidth}\raggedright
0.82\strut
\end{minipage} & \begin{minipage}[t]{0.15\columnwidth}\raggedright
1.00\strut
\end{minipage}\tabularnewline
\begin{minipage}[t]{0.13\columnwidth}\raggedright
0.15\strut
\end{minipage} & \begin{minipage}[t]{0.16\columnwidth}\raggedright
2.20\strut
\end{minipage} & \begin{minipage}[t]{0.16\columnwidth}\raggedright
2.00\strut
\end{minipage} & \begin{minipage}[t]{0.12\columnwidth}\raggedright
0.144\strut
\end{minipage} & \begin{minipage}[t]{0.10\columnwidth}\raggedright
1.27\strut
\end{minipage} & \begin{minipage}[t]{0.15\columnwidth}\raggedright
1.01\strut
\end{minipage}\tabularnewline
\begin{minipage}[t]{0.13\columnwidth}\raggedright
0.17\strut
\end{minipage} & \begin{minipage}[t]{0.16\columnwidth}\raggedright
2.18\strut
\end{minipage} & \begin{minipage}[t]{0.16\columnwidth}\raggedright
2.00\strut
\end{minipage} & \begin{minipage}[t]{0.12\columnwidth}\raggedright
0.175\strut
\end{minipage} & \begin{minipage}[t]{0.10\columnwidth}\raggedright
1.20\strut
\end{minipage} & \begin{minipage}[t]{0.15\columnwidth}\raggedright
0.98\strut
\end{minipage}\tabularnewline
\begin{minipage}[t]{0.13\columnwidth}\raggedright
0.19\strut
\end{minipage} & \begin{minipage}[t]{0.16\columnwidth}\raggedright
2.19\strut
\end{minipage} & \begin{minipage}[t]{0.16\columnwidth}\raggedright
2.00\strut
\end{minipage} & \begin{minipage}[t]{0.12\columnwidth}\raggedright
0.178\strut
\end{minipage} & \begin{minipage}[t]{0.10\columnwidth}\raggedright
0.82\strut
\end{minipage} & \begin{minipage}[t]{0.15\columnwidth}\raggedright
0.96\strut
\end{minipage}\tabularnewline
\begin{minipage}[t]{0.13\columnwidth}\raggedright
0.21\strut
\end{minipage} & \begin{minipage}[t]{0.16\columnwidth}\raggedright
2.25\strut
\end{minipage} & \begin{minipage}[t]{0.16\columnwidth}\raggedright
2.00\strut
\end{minipage} & \begin{minipage}[t]{0.12\columnwidth}\raggedright
0.173\strut
\end{minipage} & \begin{minipage}[t]{0.10\columnwidth}\raggedright
1.04\strut
\end{minipage} & \begin{minipage}[t]{0.15\columnwidth}\raggedright
1.04\strut
\end{minipage}\tabularnewline
\begin{minipage}[t]{0.13\columnwidth}\raggedright
0.23\strut
\end{minipage} & \begin{minipage}[t]{0.16\columnwidth}\raggedright
2.24\strut
\end{minipage} & \begin{minipage}[t]{0.16\columnwidth}\raggedright
2.00\strut
\end{minipage} & \begin{minipage}[t]{0.12\columnwidth}\raggedright
0.201\strut
\end{minipage} & \begin{minipage}[t]{0.10\columnwidth}\raggedright
1.34\strut
\end{minipage} & \begin{minipage}[t]{0.15\columnwidth}\raggedright
1.12\strut
\end{minipage}\tabularnewline
\begin{minipage}[t]{0.13\columnwidth}\raggedright
0.25\strut
\end{minipage} & \begin{minipage}[t]{0.16\columnwidth}\raggedright
2.28\strut
\end{minipage} & \begin{minipage}[t]{0.16\columnwidth}\raggedright
2.00\strut
\end{minipage} & \begin{minipage}[t]{0.12\columnwidth}\raggedright
0.200\strut
\end{minipage} & \begin{minipage}[t]{0.10\columnwidth}\raggedright
1.43\strut
\end{minipage} & \begin{minipage}[t]{0.15\columnwidth}\raggedright
1.14\strut
\end{minipage}\tabularnewline
\begin{minipage}[t]{0.13\columnwidth}\raggedright
0.27\strut
\end{minipage} & \begin{minipage}[t]{0.16\columnwidth}\raggedright
2.16\strut
\end{minipage} & \begin{minipage}[t]{0.16\columnwidth}\raggedright
2.00\strut
\end{minipage} & \begin{minipage}[t]{0.12\columnwidth}\raggedright
0.216\strut
\end{minipage} & \begin{minipage}[t]{0.10\columnwidth}\raggedright
1.08\strut
\end{minipage} & \begin{minipage}[t]{0.15\columnwidth}\raggedright
1.08\strut
\end{minipage}\tabularnewline
\begin{minipage}[t]{0.13\columnwidth}\raggedright
0.29\strut
\end{minipage} & \begin{minipage}[t]{0.16\columnwidth}\raggedright
1.93\strut
\end{minipage} & \begin{minipage}[t]{0.16\columnwidth}\raggedright
2.00\strut
\end{minipage} & \begin{minipage}[t]{0.12\columnwidth}\raggedright
0.241\strut
\end{minipage} & \begin{minipage}[t]{0.10\columnwidth}\raggedright
1.45\strut
\end{minipage} & \begin{minipage}[t]{0.15\columnwidth}\raggedright
0.96\strut
\end{minipage}\tabularnewline
\bottomrule
\end{longtable}

\hypertarget{appendix-b-covariance-and-correlation-matrices}{%
\subsubsection{Appendix B: Covariance and Correlation
Matrices}\label{appendix-b-covariance-and-correlation-matrices}}

The full \(13 \times 13\) total covariance matrix for the fit range
\(\tau \in [0.05, 0.30]\) is given in Table B.1 (BBB excluded from
budget). Values are in units of \([(1/N)\,dN/d\tau]^2\).
Machine-readable form: \texttt{results/covariance\_total\_fitrange.csv}.

\textbf{Table B.1: Total covariance matrix \(C_{ij}^\text{total}\) (fit
range, \(\tau \in [0.05, 0.30]\), 13 bins, BBB excluded).}

(Values \(\times 10^{-4}\) for legibility)

\begin{longtable}[]{@{}llllllllllllll@{}}
\toprule
& \(\tau\)=0.05 & 0.07 & 0.09 & 0.11 & 0.13 & 0.15 & 0.17 & 0.19 & 0.21
& 0.23 & 0.25 & 0.27 & 0.29\tabularnewline
\midrule
\endhead
\textbf{0.05} & 319.4 & 178.9 & 115.6 & 77.75 & 57.85 & 42.85 & 33.00 &
26.15 & 20.42 & 15.83 & 12.27 & 9.638 & 6.627\tabularnewline
\textbf{0.07} & 178.9 & 111.2 & 68.94 & 43.63 & 31.80 & 23.42 & 18.10 &
14.44 & 11.27 & 8.742 & 6.819 & 5.393 & 3.733\tabularnewline
\textbf{0.09} & 115.6 & 68.94 & 45.81 & 29.78 & 21.14 & 16.17 & 12.28 &
9.523 & 7.792 & 6.031 & 4.638 & 3.478 & 2.823\tabularnewline
\textbf{0.11} & 77.75 & 43.63 & 29.78 & 22.55 & 15.33 & 12.16 & 9.070 &
6.841 & 5.858 & 4.526 & 3.435 & 2.443 & 2.292\tabularnewline
\textbf{0.13} & 57.85 & 31.80 & 21.14 & 15.33 & 13.38 & 8.609 & 6.904 &
5.857 & 4.089 & 3.190 & 2.602 & 2.339 & 1.114\tabularnewline
\textbf{0.15} & 42.85 & 23.42 & 16.17 & 12.16 & 8.609 & 7.549 & 5.195 &
3.873 & 3.347 & 2.583 & 1.954 & 1.375 & 1.343\tabularnewline
\textbf{0.17} & 33.00 & 18.10 & 12.28 & 9.070 & 6.904 & 5.195 & 4.521 &
3.213 & 2.519 & 1.946 & 1.521 & 1.205 & 0.879\tabularnewline
\textbf{0.19} & 26.15 & 14.44 & 9.523 & 6.841 & 5.857 & 3.873 & 3.213 &
3.128 & 1.918 & 1.492 & 1.227 & 1.140 & 0.492\tabularnewline
\textbf{0.21} & 20.42 & 11.27 & 7.792 & 5.858 & 4.089 & 3.347 & 2.519 &
1.918 & 1.969 & 1.325 & 0.999 & 0.695 & 0.693\tabularnewline
\textbf{0.23} & 15.83 & 8.742 & 6.031 & 4.526 & 3.190 & 2.583 & 1.946 &
1.492 & 1.325 & 1.264 & 0.801 & 0.566 & 0.538\tabularnewline
\textbf{0.25} & 12.27 & 6.819 & 4.638 & 3.435 & 2.602 & 1.954 & 1.521 &
1.227 & 0.999 & 0.801 & 0.795 & 0.506 & 0.373\tabularnewline
\textbf{0.27} & 9.638 & 5.393 & 3.478 & 2.443 & 2.339 & 1.375 & 1.205 &
1.140 & 0.695 & 0.566 & 0.506 & 0.634 & 0.147\tabularnewline
\textbf{0.29} & 6.627 & 3.733 & 2.823 & 2.292 & 1.114 & 1.343 & 0.879 &
0.492 & 0.693 & 0.538 & 0.373 & 0.147 & 0.537\tabularnewline
\bottomrule
\end{longtable}

The condition number is \(3.77 \times 10^3\) (threshold per conventions:
\(10^{10}\)), reflecting well-conditioned numerics. Positive
off-diagonal elements arise primarily from the track smearing and
hadronization systematics, which are both fully correlated across bins.

\textbf{Table B.2: Correlation matrix
\(\rho_{ij} = C_{ij}/\sqrt{C_{ii} C_{jj}}\) (fit range, BBB excluded).}

\begin{longtable}[]{@{}llllllllllllll@{}}
\toprule
& 0.05 & 0.07 & 0.09 & 0.11 & 0.13 & 0.15 & 0.17 & 0.19 & 0.21 & 0.23 &
0.25 & 0.27 & 0.29\tabularnewline
\midrule
\endhead
0.05 & 1.000 & 0.949 & 0.955 & 0.916 & 0.885 & 0.873 & 0.868 & 0.827 &
0.814 & 0.788 & 0.770 & 0.677 & 0.506\tabularnewline
0.07 & 0.949 & 1.000 & 0.966 & 0.871 & 0.824 & 0.808 & 0.807 & 0.774 &
0.761 & 0.737 & 0.725 & 0.642 & 0.483\tabularnewline
0.09 & 0.955 & 0.966 & 1.000 & 0.927 & 0.854 & 0.869 & 0.853 & 0.796 &
0.820 & 0.793 & 0.768 & 0.645 & 0.569\tabularnewline
\bottomrule
\end{longtable}

(Continued rows: correlations range from 0.15 to 0.99 across the fit
range. The reduced maximum correlation compared to the BBB-inclusive
case reflects the removal of the near-rank-1 contribution.)

With BBB excluded, the correlation matrix no longer exhibits near-unity
correlations everywhere. The effective rank of the covariance matrix is
substantially higher than in the BBB-inclusive case, meaning the
\(\chi^2\) calculation uses more independent degrees of information.

\hypertarget{appendix-c-extended-cutflow-tables}{%
\subsubsection{Appendix C: Extended Cutflow
Tables}\label{appendix-c-extended-cutflow-tables}}

\textbf{Table C.1: Selection efficiency by sample and cut level.}

\begin{longtable}[]{@{}llllllll@{}}
\toprule
\begin{minipage}[b]{0.01\columnwidth}\raggedright
\strut
\end{minipage} & \begin{minipage}[b]{0.09\columnwidth}\raggedright
Data total\strut
\end{minipage} & \begin{minipage}[b]{0.26\columnwidth}\raggedright
Data \(|\cos\theta_\text{sph}|\)\strut
\end{minipage} & \begin{minipage}[b]{0.09\columnwidth}\raggedright
Data MissP\strut
\end{minipage} & \begin{minipage}[b]{0.08\columnwidth}\raggedright
Data ISR\strut
\end{minipage} & \begin{minipage}[b]{0.08\columnwidth}\raggedright
Data final\strut
\end{minipage} & \begin{minipage}[b]{0.08\columnwidth}\raggedright
MC final\strut
\end{minipage} & \begin{minipage}[b]{0.08\columnwidth}\raggedright
Data/MC\strut
\end{minipage}\tabularnewline
\midrule
\endhead
\begin{minipage}[t]{0.01\columnwidth}\raggedright
All years\strut
\end{minipage} & \begin{minipage}[t]{0.09\columnwidth}\raggedright
3,050,610\strut
\end{minipage} & \begin{minipage}[t]{0.26\columnwidth}\raggedright
97.7\%\strut
\end{minipage} & \begin{minipage}[t]{0.09\columnwidth}\raggedright
97.2\%\strut
\end{minipage} & \begin{minipage}[t]{0.08\columnwidth}\raggedright
99.0\%\strut
\end{minipage} & \begin{minipage}[t]{0.08\columnwidth}\raggedright
\textbf{94.7\%}\strut
\end{minipage} & \begin{minipage}[t]{0.08\columnwidth}\raggedright
\textbf{94.7\%}\strut
\end{minipage} & \begin{minipage}[t]{0.08\columnwidth}\raggedright
\$\textless{}\$0.1\%\strut
\end{minipage}\tabularnewline
\begin{minipage}[t]{0.01\columnwidth}\raggedright
1992\strut
\end{minipage} & \begin{minipage}[t]{0.09\columnwidth}\raggedright
551,474\strut
\end{minipage} & \begin{minipage}[t]{0.26\columnwidth}\raggedright
97.7\%\strut
\end{minipage} & \begin{minipage}[t]{0.09\columnwidth}\raggedright
97.2\%\strut
\end{minipage} & \begin{minipage}[t]{0.08\columnwidth}\raggedright
99.0\%\strut
\end{minipage} & \begin{minipage}[t]{0.08\columnwidth}\raggedright
94.7\%\strut
\end{minipage} & \begin{minipage}[t]{0.08\columnwidth}\raggedright
---\strut
\end{minipage} & \begin{minipage}[t]{0.08\columnwidth}\raggedright
---\strut
\end{minipage}\tabularnewline
\begin{minipage}[t]{0.01\columnwidth}\raggedright
1993\strut
\end{minipage} & \begin{minipage}[t]{0.09\columnwidth}\raggedright
538,601\strut
\end{minipage} & \begin{minipage}[t]{0.26\columnwidth}\raggedright
97.7\%\strut
\end{minipage} & \begin{minipage}[t]{0.09\columnwidth}\raggedright
97.2\%\strut
\end{minipage} & \begin{minipage}[t]{0.08\columnwidth}\raggedright
99.0\%\strut
\end{minipage} & \begin{minipage}[t]{0.08\columnwidth}\raggedright
94.7\%\strut
\end{minipage} & \begin{minipage}[t]{0.08\columnwidth}\raggedright
---\strut
\end{minipage} & \begin{minipage}[t]{0.08\columnwidth}\raggedright
---\strut
\end{minipage}\tabularnewline
\begin{minipage}[t]{0.01\columnwidth}\raggedright
1994 P1\strut
\end{minipage} & \begin{minipage}[t]{0.09\columnwidth}\raggedright
433,947\strut
\end{minipage} & \begin{minipage}[t]{0.26\columnwidth}\raggedright
97.7\%\strut
\end{minipage} & \begin{minipage}[t]{0.09\columnwidth}\raggedright
97.2\%\strut
\end{minipage} & \begin{minipage}[t]{0.08\columnwidth}\raggedright
99.0\%\strut
\end{minipage} & \begin{minipage}[t]{0.08\columnwidth}\raggedright
94.7\%\strut
\end{minipage} & \begin{minipage}[t]{0.08\columnwidth}\raggedright
---\strut
\end{minipage} & \begin{minipage}[t]{0.08\columnwidth}\raggedright
---\strut
\end{minipage}\tabularnewline
\begin{minipage}[t]{0.01\columnwidth}\raggedright
1994 P2\strut
\end{minipage} & \begin{minipage}[t]{0.09\columnwidth}\raggedright
447,844\strut
\end{minipage} & \begin{minipage}[t]{0.26\columnwidth}\raggedright
97.7\%\strut
\end{minipage} & \begin{minipage}[t]{0.09\columnwidth}\raggedright
97.2\%\strut
\end{minipage} & \begin{minipage}[t]{0.08\columnwidth}\raggedright
99.0\%\strut
\end{minipage} & \begin{minipage}[t]{0.08\columnwidth}\raggedright
94.7\%\strut
\end{minipage} & \begin{minipage}[t]{0.08\columnwidth}\raggedright
---\strut
\end{minipage} & \begin{minipage}[t]{0.08\columnwidth}\raggedright
---\strut
\end{minipage}\tabularnewline
\begin{minipage}[t]{0.01\columnwidth}\raggedright
1994 P3\strut
\end{minipage} & \begin{minipage}[t]{0.09\columnwidth}\raggedright
483,649\strut
\end{minipage} & \begin{minipage}[t]{0.26\columnwidth}\raggedright
97.7\%\strut
\end{minipage} & \begin{minipage}[t]{0.09\columnwidth}\raggedright
97.2\%\strut
\end{minipage} & \begin{minipage}[t]{0.08\columnwidth}\raggedright
99.0\%\strut
\end{minipage} & \begin{minipage}[t]{0.08\columnwidth}\raggedright
94.7\%\strut
\end{minipage} & \begin{minipage}[t]{0.08\columnwidth}\raggedright
---\strut
\end{minipage} & \begin{minipage}[t]{0.08\columnwidth}\raggedright
---\strut
\end{minipage}\tabularnewline
\begin{minipage}[t]{0.01\columnwidth}\raggedright
1995\strut
\end{minipage} & \begin{minipage}[t]{0.09\columnwidth}\raggedright
595,095\strut
\end{minipage} & \begin{minipage}[t]{0.26\columnwidth}\raggedright
97.7\%\strut
\end{minipage} & \begin{minipage}[t]{0.09\columnwidth}\raggedright
97.2\%\strut
\end{minipage} & \begin{minipage}[t]{0.08\columnwidth}\raggedright
99.0\%\strut
\end{minipage} & \begin{minipage}[t]{0.08\columnwidth}\raggedright
94.7\%\strut
\end{minipage} & \begin{minipage}[t]{0.08\columnwidth}\raggedright
---\strut
\end{minipage} & \begin{minipage}[t]{0.08\columnwidth}\raggedright
---\strut
\end{minipage}\tabularnewline
\bottomrule
\end{longtable}

The uniform efficiency across all years confirms stable detector
performance throughout the 1992--1995 data-taking period.

\hypertarget{appendix-d-response-matrix-properties}{%
\subsubsection{Appendix D: Response Matrix
Properties}\label{appendix-d-response-matrix-properties}}

\textbf{Table D.1: Response matrix summary.}

\begin{longtable}[]{@{}ll@{}}
\toprule
\begin{minipage}[b]{0.55\columnwidth}\raggedright
Property\strut
\end{minipage} & \begin{minipage}[b]{0.39\columnwidth}\raggedright
Value\strut
\end{minipage}\tabularnewline
\midrule
\endhead
\begin{minipage}[t]{0.55\columnwidth}\raggedright
Dimensions\strut
\end{minipage} & \begin{minipage}[t]{0.39\columnwidth}\raggedright
25 × 25\strut
\end{minipage}\tabularnewline
\begin{minipage}[t]{0.55\columnwidth}\raggedright
Bin width \(\Delta\tau\)\strut
\end{minipage} & \begin{minipage}[t]{0.39\columnwidth}\raggedright
0.02\strut
\end{minipage}\tabularnewline
\begin{minipage}[t]{0.55\columnwidth}\raggedright
\(\tau\) range\strut
\end{minipage} & \begin{minipage}[t]{0.39\columnwidth}\raggedright
\([0.00, 0.50]\)\strut
\end{minipage}\tabularnewline
\begin{minipage}[t]{0.55\columnwidth}\raggedright
Events (matched pairs)\strut
\end{minipage} & \begin{minipage}[t]{0.39\columnwidth}\raggedright
731,006\strut
\end{minipage}\tabularnewline
\begin{minipage}[t]{0.55\columnwidth}\raggedright
Column normalization\strut
\end{minipage} & \begin{minipage}[t]{0.39\columnwidth}\raggedright
1.0000 (all active bins)\strut
\end{minipage}\tabularnewline
\begin{minipage}[t]{0.55\columnwidth}\raggedright
tgenBefore events\strut
\end{minipage} & \begin{minipage}[t]{0.39\columnwidth}\raggedright
973,769\strut
\end{minipage}\tabularnewline
\begin{minipage}[t]{0.55\columnwidth}\raggedright
Generator-level efficiency \(\varepsilon_\text{gen}\)\strut
\end{minipage} & \begin{minipage}[t]{0.39\columnwidth}\raggedright
78.6\% (file 001 estimate)\strut
\end{minipage}\tabularnewline
\begin{minipage}[t]{0.55\columnwidth}\raggedright
Mean \(\tau\) bias (reco \(-\) gen)\strut
\end{minipage} & \begin{minipage}[t]{0.39\columnwidth}\raggedright
\(-0.0067\)\strut
\end{minipage}\tabularnewline
\begin{minipage}[t]{0.55\columnwidth}\raggedright
RMS smearing\strut
\end{minipage} & \begin{minipage}[t]{0.39\columnwidth}\raggedright
0.013\strut
\end{minipage}\tabularnewline
\bottomrule
\end{longtable}

\textbf{Table D.2: Response matrix diagonal fraction by bin.}

\begin{longtable}[]{@{}ll@{}}
\toprule
Bin \(\tau \in\) & Diagonal fraction\tabularnewline
\midrule
\endhead
\([0.00, 0.02]\) & 89\%\tabularnewline
\([0.02, 0.04]\) & 63\%\tabularnewline
\([0.04, 0.06]\) & 53\%\tabularnewline
\([0.06, 0.08]\) & 50\%\tabularnewline
\([0.08, 0.10]\) & 46\%\tabularnewline
\([0.10, 0.12]\) & 43\%\tabularnewline
\([0.12, 0.14]\) & 40\%\tabularnewline
\([0.14, 0.16]\) & 38\%\tabularnewline
\([0.16, 0.18]\) & 35\%\tabularnewline
\([0.18, 0.20]\) & 33\%\tabularnewline
\([0.20, 0.22]\) & 32\%\tabularnewline
\([0.22, 0.24]\) & 31\%\tabularnewline
\([0.24, 0.26]\) & 30\%\tabularnewline
\([0.26, 0.28]\) & 29\%\tabularnewline
\([0.28, 0.30]\) & 29\%\tabularnewline
\([0.30, 0.40]\) & 23--29\%\tabularnewline
\([0.40, 0.50]\) & \textasciitilde{}0\%\tabularnewline
\bottomrule
\end{longtable}

The diagonal fraction below 50\% for \(\tau > 0.04\) and below 35\% for
\(\tau > 0.10\) confirms that bin-by-bin correction is unreliable in the
fit range and IBU is essential as the primary unfolding method.

\begin{center}\rule{0.5\linewidth}{\linethickness}\end{center}

\hypertarget{references}{%
\subsection{References}\label{references}}

\begin{enumerate}
\def\labelenumi{\arabic{enumi}.}
\item
  ALEPH Collaboration, ``Studies of QCD at \(e^+e^-\) Centre-of-Mass
  Energies between 91 and 209 GeV,'' Eur. Phys. J. \textbf{C35}:457--486
  (2004); hep-ex/0409098.
\item
  S. Bethke et al. (LEP QCD Working Group), ``Combination of
  Measurements of the Strong Coupling Constant at LEP,'' hep-ex/0411006
  (2004).
\item
  ALEPH archived data analysis of thrust and two-particle correlations,
  Inspire:1793969.
\item
  G. D'Agostini, ``A multidimensional unfolding method based on Bayes'
  theorem,'' Nucl. Instrum. Methods \textbf{A362}:487--498 (1995).
\item
  T. Sjöstrand, ``PYTHIA 6.1 Physics and Manual,'' hep-ph/0010017
  (2001).
\item
  G. Dissertori, ``DISASTER++: A program for computing the NLO
  corrections to event shapes in \(e^+e^-\) annihilation,'' (1999); also
  G. Dissertori and M. Schmelling, Phys. Lett. \textbf{B361}:167--178
  (1995).
\item
  R. Abbate, M. Fickinger, A. H. Hoang, V. Mateu, I. W. Stewart,
  ``Thrust at N\(^3\)LL with Power Corrections and a Precision Global
  Fit for \(\alpha_s(m_Z)\),'' Phys. Rev. \textbf{D83}:074021 (2011);
  arXiv:1006.3080.
\item
  ALEPH Collaboration, ``Performance of the ALEPH detector at LEP,''
  Nucl. Instrum. Methods \textbf{A360}:481--506 (1995).
\item
  Particle Data Group, ``Review of Particle Physics,'' Prog. Theor. Exp.
  Phys. \textbf{2022}:083C01 (2022).
\item
  V. A. Khoze, W. Ochs, J. Wosiek, ``Particle production mechanisms,''
  hep-ph/0009298 (2000) --- review of Lund string model and
  fragmentation.
\end{enumerate}

\begin{center}\rule{0.5\linewidth}{\linethickness}\end{center}

\hypertarget{machine-readable-results}{%
\subsection{Machine-Readable Results}\label{machine-readable-results}}

The following files are provided in the \texttt{results/} subdirectory:

\begin{itemize}
\tightlist
\item
  \texttt{results/thrust\_distribution.csv}: Full corrected thrust
  distribution with per-bin statistical and systematic uncertainties (13
  bins in fit range).
\item
  \texttt{results/covariance\_total\_fitrange.csv}: \(13 \times 13\)
  total covariance matrix for the fit range \(\tau \in [0.05, 0.30]\).
\item
  \texttt{results/alphas\_result.csv}: Indicative \(\alpha_s\) result
  and scale variation envelope.
\end{itemize}

Additional machine-readable files from Phase 4 (not copied to this
directory but available in \texttt{../../phase4\_inference/exec/}):

\begin{itemize}
\tightlist
\item
  \texttt{covariance\_stat.npz}: Statistical covariance matrix (25×25)
\item
  \texttt{covariance\_syst.npz}: Systematic covariance matrix (25×25)
\item
  \texttt{covariance\_total.npz}: Total covariance matrix (25×25)
\item
  \texttt{systematics\_shifts.npz}: Per-source systematic shifts
\item
  \texttt{indep\_closure\_results.npz}: Independent closure test results
\item
  \texttt{stress\_test\_results.npz}: Stress test results
\end{itemize}

\begin{center}\rule{0.5\linewidth}{\linethickness}\end{center}

\emph{Analysis completed: Phase 5. All scripts available in
\texttt{phase3\_selection/scripts/} and
\texttt{phase4\_inference/scripts/}. Full reproducible analysis chain:
\texttt{pixi\ run\ all} from the analysis directory
\texttt{thrust\_alphas/}.}

\end{document}
